\documentclass[11pt]{article}

    \usepackage[breakable]{tcolorbox}
    \usepackage{parskip} % Stop auto-indenting (to mimic markdown behaviour)
    

    % Basic figure setup, for now with no caption control since it's done
    % automatically by Pandoc (which extracts ![](path) syntax from Markdown).
    \usepackage{graphicx}
    % Keep aspect ratio if custom image width or height is specified
    \setkeys{Gin}{keepaspectratio}
    % Maintain compatibility with old templates. Remove in nbconvert 6.0
    \let\Oldincludegraphics\includegraphics
    % Ensure that by default, figures have no caption (until we provide a
    % proper Figure object with a Caption API and a way to capture that
    % in the conversion process - todo).
    \usepackage{caption}
    \DeclareCaptionFormat{nocaption}{}
    \captionsetup{format=nocaption,aboveskip=0pt,belowskip=0pt}

    \usepackage{float}
    \floatplacement{figure}{H} % forces figures to be placed at the correct location
    \usepackage{xcolor} % Allow colors to be defined
    \usepackage{enumerate} % Needed for markdown enumerations to work
    \usepackage{geometry} % Used to adjust the document margins
    \usepackage{amsmath} % Equations
    \usepackage{amssymb} % Equations
    \usepackage{textcomp} % defines textquotesingle
    % Hack from http://tex.stackexchange.com/a/47451/13684:
    \AtBeginDocument{%
        \def\PYZsq{\textquotesingle}% Upright quotes in Pygmentized code
    }
    \usepackage{upquote} % Upright quotes for verbatim code
    \usepackage{eurosym} % defines \euro

    \usepackage{iftex}
    \ifPDFTeX
        \usepackage[T1]{fontenc}
        \IfFileExists{alphabeta.sty}{
              \usepackage{alphabeta}
          }{
              \usepackage[utf8]{inputenc}
              \DeclareUnicodeCharacter{03B2}{\ensuremath{\beta}}
          }
    \else
        \usepackage{fontspec}
        \usepackage{unicode-math}
    \fi

    \usepackage{fancyvrb} % verbatim replacement that allows latex
    \usepackage{grffile} % extends the file name processing of package graphics
                         % to support a larger range
    \makeatletter % fix for old versions of grffile with XeLaTeX
    \@ifpackagelater{grffile}{2019/11/01}
    {
      % Do nothing on new versions
    }
    {
      \def\Gread@@xetex#1{%
        \IfFileExists{"\Gin@base".bb}%
        {\Gread@eps{\Gin@base.bb}}%
        {\Gread@@xetex@aux#1}%
      }
    }
    \makeatother
    \usepackage[Export]{adjustbox} % Used to constrain images to a maximum size
    \adjustboxset{max size={0.9\linewidth}{0.9\paperheight}}

    % The hyperref package gives us a pdf with properly built
    % internal navigation ('pdf bookmarks' for the table of contents,
    % internal cross-reference links, web links for URLs, etc.)
    \usepackage{hyperref}
    % The default LaTeX title has an obnoxious amount of whitespace. By default,
    % titling removes some of it. It also provides customization options.
    \usepackage{titling}
    \usepackage{longtable} % longtable support required by pandoc >1.10
    \usepackage{booktabs}  % table support for pandoc > 1.12.2
    \usepackage{array}     % table support for pandoc >= 2.11.3
    \usepackage{calc}      % table minipage width calculation for pandoc >= 2.11.1
    \usepackage[inline]{enumitem} % IRkernel/repr support (it uses the enumerate* environment)
    \usepackage[normalem]{ulem} % ulem is needed to support strikethroughs (\sout)
                                % normalem makes italics be italics, not underlines
    \usepackage{soul}      % strikethrough (\st) support for pandoc >= 3.0.0
    \usepackage{mathrsfs}
    

    
    % Colors for the hyperref package
    \definecolor{urlcolor}{rgb}{0,.145,.698}
    \definecolor{linkcolor}{rgb}{.71,0.21,0.01}
    \definecolor{citecolor}{rgb}{.12,.54,.11}

    % ANSI colors
    \definecolor{ansi-black}{HTML}{3E424D}
    \definecolor{ansi-black-intense}{HTML}{282C36}
    \definecolor{ansi-red}{HTML}{E75C58}
    \definecolor{ansi-red-intense}{HTML}{B22B31}
    \definecolor{ansi-green}{HTML}{00A250}
    \definecolor{ansi-green-intense}{HTML}{007427}
    \definecolor{ansi-yellow}{HTML}{DDB62B}
    \definecolor{ansi-yellow-intense}{HTML}{B27D12}
    \definecolor{ansi-blue}{HTML}{208FFB}
    \definecolor{ansi-blue-intense}{HTML}{0065CA}
    \definecolor{ansi-magenta}{HTML}{D160C4}
    \definecolor{ansi-magenta-intense}{HTML}{A03196}
    \definecolor{ansi-cyan}{HTML}{60C6C8}
    \definecolor{ansi-cyan-intense}{HTML}{258F8F}
    \definecolor{ansi-white}{HTML}{C5C1B4}
    \definecolor{ansi-white-intense}{HTML}{A1A6B2}
    \definecolor{ansi-default-inverse-fg}{HTML}{FFFFFF}
    \definecolor{ansi-default-inverse-bg}{HTML}{000000}

    % common color for the border for error outputs.
    \definecolor{outerrorbackground}{HTML}{FFDFDF}

    % commands and environments needed by pandoc snippets
    % extracted from the output of `pandoc -s`
    \providecommand{\tightlist}{%
      \setlength{\itemsep}{0pt}\setlength{\parskip}{0pt}}
    \DefineVerbatimEnvironment{Highlighting}{Verbatim}{commandchars=\\\{\}}
    % Add ',fontsize=\small' for more characters per line
    \newenvironment{Shaded}{}{}
    \newcommand{\KeywordTok}[1]{\textcolor[rgb]{0.00,0.44,0.13}{\textbf{{#1}}}}
    \newcommand{\DataTypeTok}[1]{\textcolor[rgb]{0.56,0.13,0.00}{{#1}}}
    \newcommand{\DecValTok}[1]{\textcolor[rgb]{0.25,0.63,0.44}{{#1}}}
    \newcommand{\BaseNTok}[1]{\textcolor[rgb]{0.25,0.63,0.44}{{#1}}}
    \newcommand{\FloatTok}[1]{\textcolor[rgb]{0.25,0.63,0.44}{{#1}}}
    \newcommand{\CharTok}[1]{\textcolor[rgb]{0.25,0.44,0.63}{{#1}}}
    \newcommand{\StringTok}[1]{\textcolor[rgb]{0.25,0.44,0.63}{{#1}}}
    \newcommand{\CommentTok}[1]{\textcolor[rgb]{0.38,0.63,0.69}{\textit{{#1}}}}
    \newcommand{\OtherTok}[1]{\textcolor[rgb]{0.00,0.44,0.13}{{#1}}}
    \newcommand{\AlertTok}[1]{\textcolor[rgb]{1.00,0.00,0.00}{\textbf{{#1}}}}
    \newcommand{\FunctionTok}[1]{\textcolor[rgb]{0.02,0.16,0.49}{{#1}}}
    \newcommand{\RegionMarkerTok}[1]{{#1}}
    \newcommand{\ErrorTok}[1]{\textcolor[rgb]{1.00,0.00,0.00}{\textbf{{#1}}}}
    \newcommand{\NormalTok}[1]{{#1}}

    % Additional commands for more recent versions of Pandoc
    \newcommand{\ConstantTok}[1]{\textcolor[rgb]{0.53,0.00,0.00}{{#1}}}
    \newcommand{\SpecialCharTok}[1]{\textcolor[rgb]{0.25,0.44,0.63}{{#1}}}
    \newcommand{\VerbatimStringTok}[1]{\textcolor[rgb]{0.25,0.44,0.63}{{#1}}}
    \newcommand{\SpecialStringTok}[1]{\textcolor[rgb]{0.73,0.40,0.53}{{#1}}}
    \newcommand{\ImportTok}[1]{{#1}}
    \newcommand{\DocumentationTok}[1]{\textcolor[rgb]{0.73,0.13,0.13}{\textit{{#1}}}}
    \newcommand{\AnnotationTok}[1]{\textcolor[rgb]{0.38,0.63,0.69}{\textbf{\textit{{#1}}}}}
    \newcommand{\CommentVarTok}[1]{\textcolor[rgb]{0.38,0.63,0.69}{\textbf{\textit{{#1}}}}}
    \newcommand{\VariableTok}[1]{\textcolor[rgb]{0.10,0.09,0.49}{{#1}}}
    \newcommand{\ControlFlowTok}[1]{\textcolor[rgb]{0.00,0.44,0.13}{\textbf{{#1}}}}
    \newcommand{\OperatorTok}[1]{\textcolor[rgb]{0.40,0.40,0.40}{{#1}}}
    \newcommand{\BuiltInTok}[1]{{#1}}
    \newcommand{\ExtensionTok}[1]{{#1}}
    \newcommand{\PreprocessorTok}[1]{\textcolor[rgb]{0.74,0.48,0.00}{{#1}}}
    \newcommand{\AttributeTok}[1]{\textcolor[rgb]{0.49,0.56,0.16}{{#1}}}
    \newcommand{\InformationTok}[1]{\textcolor[rgb]{0.38,0.63,0.69}{\textbf{\textit{{#1}}}}}
    \newcommand{\WarningTok}[1]{\textcolor[rgb]{0.38,0.63,0.69}{\textbf{\textit{{#1}}}}}
    \makeatletter
    \newsavebox\pandoc@box
    \newcommand*\pandocbounded[1]{%
      \sbox\pandoc@box{#1}%
      % scaling factors for width and height
      \Gscale@div\@tempa\textheight{\dimexpr\ht\pandoc@box+\dp\pandoc@box\relax}%
      \Gscale@div\@tempb\linewidth{\wd\pandoc@box}%
      % select the smaller of both
      \ifdim\@tempb\p@<\@tempa\p@
        \let\@tempa\@tempb
      \fi
      % scaling accordingly (\@tempa < 1)
      \ifdim\@tempa\p@<\p@
        \scalebox{\@tempa}{\usebox\pandoc@box}%
      % scaling not needed, use as it is
      \else
        \usebox{\pandoc@box}%
      \fi
    }
    \makeatother

    % Define a nice break command that doesn't care if a line doesn't already
    % exist.
    \def\br{\hspace*{\fill} \\* }
    % Math Jax compatibility definitions
    \def\gt{>}
    \def\lt{<}
    \let\Oldtex\TeX
    \let\Oldlatex\LaTeX
    \renewcommand{\TeX}{\textrm{\Oldtex}}
    \renewcommand{\LaTeX}{\textrm{\Oldlatex}}
    % Document parameters
    % Document title
    \title{Homework 4\\[0.5ex]\large Computer Simulation 1}
    \author{Andrea Pokraj\v{c}i\'{c}\\
    \normalsize\href{mailto:andrea.pokrajcic@studserv.uni-leipzig.de}{andrea.pokrajcic@studserv.uni-leipzig.de}\\
    \normalsize Matriculation number: 3793956}
    \date{}
    
    
    
    
    
    
    
% Pygments definitions
\makeatletter
\def\PY@reset{\let\PY@it=\relax \let\PY@bf=\relax%
    \let\PY@ul=\relax \let\PY@tc=\relax%
    \let\PY@bc=\relax \let\PY@ff=\relax}
\def\PY@tok#1{\csname PY@tok@#1\endcsname}
\def\PY@toks#1+{\ifx\relax#1\empty\else%
    \PY@tok{#1}\expandafter\PY@toks\fi}
\def\PY@do#1{\PY@bc{\PY@tc{\PY@ul{%
    \PY@it{\PY@bf{\PY@ff{#1}}}}}}}
\def\PY#1#2{\PY@reset\PY@toks#1+\relax+\PY@do{#2}}

\@namedef{PY@tok@w}{\def\PY@tc##1{\textcolor[rgb]{0.73,0.73,0.73}{##1}}}
\@namedef{PY@tok@c}{\let\PY@it=\textit\def\PY@tc##1{\textcolor[rgb]{0.24,0.48,0.48}{##1}}}
\@namedef{PY@tok@cp}{\def\PY@tc##1{\textcolor[rgb]{0.61,0.40,0.00}{##1}}}
\@namedef{PY@tok@k}{\let\PY@bf=\textbf\def\PY@tc##1{\textcolor[rgb]{0.00,0.50,0.00}{##1}}}
\@namedef{PY@tok@kp}{\def\PY@tc##1{\textcolor[rgb]{0.00,0.50,0.00}{##1}}}
\@namedef{PY@tok@kt}{\def\PY@tc##1{\textcolor[rgb]{0.69,0.00,0.25}{##1}}}
\@namedef{PY@tok@o}{\def\PY@tc##1{\textcolor[rgb]{0.40,0.40,0.40}{##1}}}
\@namedef{PY@tok@ow}{\let\PY@bf=\textbf\def\PY@tc##1{\textcolor[rgb]{0.67,0.13,1.00}{##1}}}
\@namedef{PY@tok@nb}{\def\PY@tc##1{\textcolor[rgb]{0.00,0.50,0.00}{##1}}}
\@namedef{PY@tok@nf}{\def\PY@tc##1{\textcolor[rgb]{0.00,0.00,1.00}{##1}}}
\@namedef{PY@tok@nc}{\let\PY@bf=\textbf\def\PY@tc##1{\textcolor[rgb]{0.00,0.00,1.00}{##1}}}
\@namedef{PY@tok@nn}{\let\PY@bf=\textbf\def\PY@tc##1{\textcolor[rgb]{0.00,0.00,1.00}{##1}}}
\@namedef{PY@tok@ne}{\let\PY@bf=\textbf\def\PY@tc##1{\textcolor[rgb]{0.80,0.25,0.22}{##1}}}
\@namedef{PY@tok@nv}{\def\PY@tc##1{\textcolor[rgb]{0.10,0.09,0.49}{##1}}}
\@namedef{PY@tok@no}{\def\PY@tc##1{\textcolor[rgb]{0.53,0.00,0.00}{##1}}}
\@namedef{PY@tok@nl}{\def\PY@tc##1{\textcolor[rgb]{0.46,0.46,0.00}{##1}}}
\@namedef{PY@tok@ni}{\let\PY@bf=\textbf\def\PY@tc##1{\textcolor[rgb]{0.44,0.44,0.44}{##1}}}
\@namedef{PY@tok@na}{\def\PY@tc##1{\textcolor[rgb]{0.41,0.47,0.13}{##1}}}
\@namedef{PY@tok@nt}{\let\PY@bf=\textbf\def\PY@tc##1{\textcolor[rgb]{0.00,0.50,0.00}{##1}}}
\@namedef{PY@tok@nd}{\def\PY@tc##1{\textcolor[rgb]{0.67,0.13,1.00}{##1}}}
\@namedef{PY@tok@s}{\def\PY@tc##1{\textcolor[rgb]{0.73,0.13,0.13}{##1}}}
\@namedef{PY@tok@sd}{\let\PY@it=\textit\def\PY@tc##1{\textcolor[rgb]{0.73,0.13,0.13}{##1}}}
\@namedef{PY@tok@si}{\let\PY@bf=\textbf\def\PY@tc##1{\textcolor[rgb]{0.64,0.35,0.47}{##1}}}
\@namedef{PY@tok@se}{\let\PY@bf=\textbf\def\PY@tc##1{\textcolor[rgb]{0.67,0.36,0.12}{##1}}}
\@namedef{PY@tok@sr}{\def\PY@tc##1{\textcolor[rgb]{0.64,0.35,0.47}{##1}}}
\@namedef{PY@tok@ss}{\def\PY@tc##1{\textcolor[rgb]{0.10,0.09,0.49}{##1}}}
\@namedef{PY@tok@sx}{\def\PY@tc##1{\textcolor[rgb]{0.00,0.50,0.00}{##1}}}
\@namedef{PY@tok@m}{\def\PY@tc##1{\textcolor[rgb]{0.40,0.40,0.40}{##1}}}
\@namedef{PY@tok@gh}{\let\PY@bf=\textbf\def\PY@tc##1{\textcolor[rgb]{0.00,0.00,0.50}{##1}}}
\@namedef{PY@tok@gu}{\let\PY@bf=\textbf\def\PY@tc##1{\textcolor[rgb]{0.50,0.00,0.50}{##1}}}
\@namedef{PY@tok@gd}{\def\PY@tc##1{\textcolor[rgb]{0.63,0.00,0.00}{##1}}}
\@namedef{PY@tok@gi}{\def\PY@tc##1{\textcolor[rgb]{0.00,0.52,0.00}{##1}}}
\@namedef{PY@tok@gr}{\def\PY@tc##1{\textcolor[rgb]{0.89,0.00,0.00}{##1}}}
\@namedef{PY@tok@ge}{\let\PY@it=\textit}
\@namedef{PY@tok@gs}{\let\PY@bf=\textbf}
\@namedef{PY@tok@ges}{\let\PY@bf=\textbf\let\PY@it=\textit}
\@namedef{PY@tok@gp}{\let\PY@bf=\textbf\def\PY@tc##1{\textcolor[rgb]{0.00,0.00,0.50}{##1}}}
\@namedef{PY@tok@go}{\def\PY@tc##1{\textcolor[rgb]{0.44,0.44,0.44}{##1}}}
\@namedef{PY@tok@gt}{\def\PY@tc##1{\textcolor[rgb]{0.00,0.27,0.87}{##1}}}
\@namedef{PY@tok@err}{\def\PY@bc##1{{\setlength{\fboxsep}{\string -\fboxrule}\fcolorbox[rgb]{1.00,0.00,0.00}{1,1,1}{\strut ##1}}}}
\@namedef{PY@tok@kc}{\let\PY@bf=\textbf\def\PY@tc##1{\textcolor[rgb]{0.00,0.50,0.00}{##1}}}
\@namedef{PY@tok@kd}{\let\PY@bf=\textbf\def\PY@tc##1{\textcolor[rgb]{0.00,0.50,0.00}{##1}}}
\@namedef{PY@tok@kn}{\let\PY@bf=\textbf\def\PY@tc##1{\textcolor[rgb]{0.00,0.50,0.00}{##1}}}
\@namedef{PY@tok@kr}{\let\PY@bf=\textbf\def\PY@tc##1{\textcolor[rgb]{0.00,0.50,0.00}{##1}}}
\@namedef{PY@tok@bp}{\def\PY@tc##1{\textcolor[rgb]{0.00,0.50,0.00}{##1}}}
\@namedef{PY@tok@fm}{\def\PY@tc##1{\textcolor[rgb]{0.00,0.00,1.00}{##1}}}
\@namedef{PY@tok@vc}{\def\PY@tc##1{\textcolor[rgb]{0.10,0.09,0.49}{##1}}}
\@namedef{PY@tok@vg}{\def\PY@tc##1{\textcolor[rgb]{0.10,0.09,0.49}{##1}}}
\@namedef{PY@tok@vi}{\def\PY@tc##1{\textcolor[rgb]{0.10,0.09,0.49}{##1}}}
\@namedef{PY@tok@vm}{\def\PY@tc##1{\textcolor[rgb]{0.10,0.09,0.49}{##1}}}
\@namedef{PY@tok@sa}{\def\PY@tc##1{\textcolor[rgb]{0.73,0.13,0.13}{##1}}}
\@namedef{PY@tok@sb}{\def\PY@tc##1{\textcolor[rgb]{0.73,0.13,0.13}{##1}}}
\@namedef{PY@tok@sc}{\def\PY@tc##1{\textcolor[rgb]{0.73,0.13,0.13}{##1}}}
\@namedef{PY@tok@dl}{\def\PY@tc##1{\textcolor[rgb]{0.73,0.13,0.13}{##1}}}
\@namedef{PY@tok@s2}{\def\PY@tc##1{\textcolor[rgb]{0.73,0.13,0.13}{##1}}}
\@namedef{PY@tok@sh}{\def\PY@tc##1{\textcolor[rgb]{0.73,0.13,0.13}{##1}}}
\@namedef{PY@tok@s1}{\def\PY@tc##1{\textcolor[rgb]{0.73,0.13,0.13}{##1}}}
\@namedef{PY@tok@mb}{\def\PY@tc##1{\textcolor[rgb]{0.40,0.40,0.40}{##1}}}
\@namedef{PY@tok@mf}{\def\PY@tc##1{\textcolor[rgb]{0.40,0.40,0.40}{##1}}}
\@namedef{PY@tok@mh}{\def\PY@tc##1{\textcolor[rgb]{0.40,0.40,0.40}{##1}}}
\@namedef{PY@tok@mi}{\def\PY@tc##1{\textcolor[rgb]{0.40,0.40,0.40}{##1}}}
\@namedef{PY@tok@il}{\def\PY@tc##1{\textcolor[rgb]{0.40,0.40,0.40}{##1}}}
\@namedef{PY@tok@mo}{\def\PY@tc##1{\textcolor[rgb]{0.40,0.40,0.40}{##1}}}
\@namedef{PY@tok@ch}{\let\PY@it=\textit\def\PY@tc##1{\textcolor[rgb]{0.24,0.48,0.48}{##1}}}
\@namedef{PY@tok@cm}{\let\PY@it=\textit\def\PY@tc##1{\textcolor[rgb]{0.24,0.48,0.48}{##1}}}
\@namedef{PY@tok@cpf}{\let\PY@it=\textit\def\PY@tc##1{\textcolor[rgb]{0.24,0.48,0.48}{##1}}}
\@namedef{PY@tok@c1}{\let\PY@it=\textit\def\PY@tc##1{\textcolor[rgb]{0.24,0.48,0.48}{##1}}}
\@namedef{PY@tok@cs}{\let\PY@it=\textit\def\PY@tc##1{\textcolor[rgb]{0.24,0.48,0.48}{##1}}}

\def\PYZbs{\char`\\}
\def\PYZus{\char`\_}
\def\PYZob{\char`\{}
\def\PYZcb{\char`\}}
\def\PYZca{\char`\^}
\def\PYZam{\char`\&}
\def\PYZlt{\char`\<}
\def\PYZgt{\char`\>}
\def\PYZsh{\char`\#}
\def\PYZpc{\char`\%}
\def\PYZdl{\char`\$}
\def\PYZhy{\char`\-}
\def\PYZsq{\char`\'}
\def\PYZdq{\char`\"}
\def\PYZti{\char`\~}
% for compatibility with earlier versions
\def\PYZat{@}
\def\PYZlb{[}
\def\PYZrb{]}
\makeatother


    % For linebreaks inside Verbatim environment from package fancyvrb.
    \makeatletter
        \newbox\Wrappedcontinuationbox
        \newbox\Wrappedvisiblespacebox
        \newcommand*\Wrappedvisiblespace {\textcolor{red}{\textvisiblespace}}
        \newcommand*\Wrappedcontinuationsymbol {\textcolor{red}{\llap{\tiny$\m@th\hookrightarrow$}}}
        \newcommand*\Wrappedcontinuationindent {3ex }
        \newcommand*\Wrappedafterbreak {\kern\Wrappedcontinuationindent\copy\Wrappedcontinuationbox}
        % Take advantage of the already applied Pygments mark-up to insert
        % potential linebreaks for TeX processing.
        %        {, <, #, %, $, ' and ": go to next line.
        %        _, }, ^, &, >, - and ~: stay at end of broken line.
        % Use of \textquotesingle for straight quote.
        \newcommand*\Wrappedbreaksatspecials {%
            \def\PYGZus{\discretionary{\char`\_}{\Wrappedafterbreak}{\char`\_}}%
            \def\PYGZob{\discretionary{}{\Wrappedafterbreak\char`\{}{\char`\{}}%
            \def\PYGZcb{\discretionary{\char`\}}{\Wrappedafterbreak}{\char`\}}}%
            \def\PYGZca{\discretionary{\char`\^}{\Wrappedafterbreak}{\char`\^}}%
            \def\PYGZam{\discretionary{\char`\&}{\Wrappedafterbreak}{\char`\&}}%
            \def\PYGZlt{\discretionary{}{\Wrappedafterbreak\char`\<}{\char`\<}}%
            \def\PYGZgt{\discretionary{\char`\>}{\Wrappedafterbreak}{\char`\>}}%
            \def\PYGZsh{\discretionary{}{\Wrappedafterbreak\char`\#}{\char`\#}}%
            \def\PYGZpc{\discretionary{}{\Wrappedafterbreak\char`\%}{\char`\%}}%
            \def\PYGZdl{\discretionary{}{\Wrappedafterbreak\char`\$}{\char`\$}}%
            \def\PYGZhy{\discretionary{\char`\-}{\Wrappedafterbreak}{\char`\-}}%
            \def\PYGZsq{\discretionary{}{\Wrappedafterbreak\textquotesingle}{\textquotesingle}}%
            \def\PYGZdq{\discretionary{}{\Wrappedafterbreak\char`\"}{\char`\"}}%
            \def\PYGZti{\discretionary{\char`\~}{\Wrappedafterbreak}{\char`\~}}%
        }
        % Some characters . , ; ? ! / are not pygmentized.
        % This macro makes them "active" and they will insert potential linebreaks
        \newcommand*\Wrappedbreaksatpunct {%
            \lccode`\~`\.\lowercase{\def~}{\discretionary{\hbox{\char`\.}}{\Wrappedafterbreak}{\hbox{\char`\.}}}%
            \lccode`\~`\,\lowercase{\def~}{\discretionary{\hbox{\char`\,}}{\Wrappedafterbreak}{\hbox{\char`\,}}}%
            \lccode`\~`\;\lowercase{\def~}{\discretionary{\hbox{\char`\;}}{\Wrappedafterbreak}{\hbox{\char`\;}}}%
            \lccode`\~`\:\lowercase{\def~}{\discretionary{\hbox{\char`\:}}{\Wrappedafterbreak}{\hbox{\char`\:}}}%
            \lccode`\~`\?\lowercase{\def~}{\discretionary{\hbox{\char`\?}}{\Wrappedafterbreak}{\hbox{\char`\?}}}%
            \lccode`\~`\!\lowercase{\def~}{\discretionary{\hbox{\char`\!}}{\Wrappedafterbreak}{\hbox{\char`\!}}}%
            \lccode`\~`\/\lowercase{\def~}{\discretionary{\hbox{\char`\/}}{\Wrappedafterbreak}{\hbox{\char`\/}}}%
            \catcode`\.\active
            \catcode`\,\active
            \catcode`\;\active
            \catcode`\:\active
            \catcode`\?\active
            \catcode`\!\active
            \catcode`\/\active
            \lccode`\~`\~
        }
    \makeatother

    \let\OriginalVerbatim=\Verbatim
    \makeatletter
    \renewcommand{\Verbatim}[1][1]{%
        %\parskip\z@skip
        \sbox\Wrappedcontinuationbox {\Wrappedcontinuationsymbol}%
        \sbox\Wrappedvisiblespacebox {\FV@SetupFont\Wrappedvisiblespace}%
        \def\FancyVerbFormatLine ##1{\hsize\linewidth
            \vtop{\raggedright\hyphenpenalty\z@\exhyphenpenalty\z@
                \doublehyphendemerits\z@\finalhyphendemerits\z@
                \strut ##1\strut}%
        }%
        % If the linebreak is at a space, the latter will be displayed as visible
        % space at end of first line, and a continuation symbol starts next line.
        % Stretch/shrink are however usually zero for typewriter font.
        \def\FV@Space {%
            \nobreak\hskip\z@ plus\fontdimen3\font minus\fontdimen4\font
            \discretionary{\copy\Wrappedvisiblespacebox}{\Wrappedafterbreak}
            {\kern\fontdimen2\font}%
        }%

        % Allow breaks at special characters using \PYG... macros.
        \Wrappedbreaksatspecials
        % Breaks at punctuation characters . , ; ? ! and / need catcode=\active
        \OriginalVerbatim[#1,codes*=\Wrappedbreaksatpunct]%
    }
    \makeatother

    % Exact colors from NB
    \definecolor{incolor}{HTML}{303F9F}
    \definecolor{outcolor}{HTML}{D84315}
    \definecolor{cellborder}{HTML}{CFCFCF}
    \definecolor{cellbackground}{HTML}{F7F7F7}

    % prompt
    \makeatletter
    \newcommand{\boxspacing}{\kern\kvtcb@left@rule\kern\kvtcb@boxsep}
    \makeatother
    \newcommand{\prompt}[4]{
        {\ttfamily\llap{{\color{#2}[#3]:\hspace{3pt}#4}}\vspace{-\baselineskip}}
    }
    

    
    % Prevent overflowing lines due to hard-to-break entities
    \sloppy
    % Setup hyperref package
    \hypersetup{
      breaklinks=true,  % so long urls are correctly broken across lines
      colorlinks=true,
      urlcolor=urlcolor,
      linkcolor=linkcolor,
      citecolor=citecolor,
      }
    % Slightly bigger margins than the latex defaults
    
    \geometry{verbose,tmargin=1in,bmargin=1in,lmargin=1in,rmargin=1in}
    
    

\begin{document}
    
    \maketitle
    
    

    
    \section*{Problem 7: Statement}\label{problem-7-statement}

The generation of equally distributed pseudo random numbers with
so-called \emph{Linear Congruential Generator} (LCG) proceeds according
to the rule

\[ X_{i+1}= (aX_i +c)\ mod \ m \]

where for the following parts of this homework the (intentionally
non-optimal) param3eters should be chosen as \(a=5\), \(c=0\),
\(m=2**11=2048\), \(X_0=1\).

\begin{enumerate}
\def\labelenumi{(\alph{enumi})}
\item
  Implement an LCG with these parameters and choose an appropriatte
  normalisation such that \(X_i \in [0,1)\).
\item
  Plot the time series \(X_i\) of \(10^4\) sucessive random numbers and
  determine their period \(\tau\). Which meaning has this quantity?
\item
  Compute the average \(\bar{X} =  \sum_{i=1}^N X_i/N\) and the
  estimator for the variance \(\hat{\sigma}_X^2=\bar{X^2}-\bar{X}^2\)
  for \(N = \tau , 10\tau, 100\tau\) random numbers. Why do the resulčts
  not improve isnpite the higher statistics?
\item
  Plot the so-called running average \(\bar{X} (n) =\sum_{i=1}^n X_i/n\)
  as function of n in the interval \(n \in [30 000, 40 000]\).
\item
  Plot an xy graph with \(x= X_i\) and \(y=X_{i+1}\) for \(i=1,3,5,...\)
\end{enumerate}

    \section*{Problem 7: Solution}\label{problem-7-solution}

    \subsection*{Part (a)}

The linear congruential generator (LCG) generates a sequence of integers
according to

\[
X_{i+1}=(aX_i+c) \bmod m \tag{1}
\]

with \(a=5\), \(c=0\), \(m=2^{11}=2048\) and \(X_0=1\). The remainder of
a division by \(m\) is always one of the integers \(0,1,\dots,m-1\).
Therefore, if we divide by \(m\) we get numbers in the interval
\([0,1)\):

\[
U_i=\frac{X_i}{m} \tag{2}
\]

We implement the generator and generate \(10^5\) numbers, which is
enough for all parts of this problem.

    \begin{tcolorbox}[breakable, size=fbox, boxrule=1pt, pad at break*=1mm,colback=cellbackground, colframe=cellborder]
\prompt{In}{incolor}{9}{\boxspacing}
\begin{Verbatim}[commandchars=\\\{\}]
\PY{k+kn}{import} \PY{n+nn}{numpy} \PY{k}{as} \PY{n+nn}{np}
\PY{k+kn}{from} \PY{n+nn}{matplotlib} \PY{k+kn}{import} \PY{n}{pyplot} \PY{k}{as} \PY{n}{plt}

\PY{n}{a} \PY{o}{=} \PY{l+m+mi}{5}
\PY{n}{c} \PY{o}{=} \PY{l+m+mi}{0}
\PY{n}{m} \PY{o}{=} \PY{l+m+mi}{2}\PY{o}{*}\PY{o}{*}\PY{l+m+mi}{11}
\PY{n}{X\PYZus{}0} \PY{o}{=} \PY{l+m+mi}{1}

\PY{k}{def} \PY{n+nf}{LCG}\PY{p}{(}\PY{n}{n}\PY{p}{)}\PY{p}{:}
    \PY{n}{X} \PY{o}{=} \PY{p}{[}\PY{n}{X\PYZus{}0}\PY{p}{]}
    \PY{k}{for} \PY{n}{i} \PY{o+ow}{in} \PY{n+nb}{range}\PY{p}{(}\PY{n}{n}\PY{p}{)}\PY{p}{:}
        \PY{n}{X}\PY{o}{.}\PY{n}{append}\PY{p}{(}\PY{p}{(}\PY{n}{a}\PY{o}{*}\PY{n}{X}\PY{p}{[}\PY{o}{\PYZhy{}}\PY{l+m+mi}{1}\PY{p}{]}\PY{o}{+}\PY{n}{c}\PY{p}{)} \PY{o}{\PYZpc{}} \PY{n}{m}\PY{p}{)}
    \PY{k}{return} \PY{n}{X}

\PY{n}{X} \PY{o}{=} \PY{n}{LCG}\PY{p}{(}\PY{l+m+mi}{10}\PY{o}{*}\PY{o}{*}\PY{l+m+mi}{5}\PY{p}{)}
\PY{n}{U} \PY{o}{=} \PY{n}{np}\PY{o}{.}\PY{n}{array}\PY{p}{(}\PY{n}{X}\PY{p}{)}\PY{o}{/}\PY{n}{m}

\PY{n+nb}{print}\PY{p}{(}\PY{l+s+s2}{\PYZdq{}}\PY{l+s+s2}{First 10 X\PYZus{}i:}\PY{l+s+s2}{\PYZdq{}}\PY{p}{,} \PY{n}{X}\PY{p}{[}\PY{p}{:}\PY{l+m+mi}{10}\PY{p}{]}\PY{p}{)}
\PY{n+nb}{print}\PY{p}{(}\PY{l+s+s2}{\PYZdq{}}\PY{l+s+s2}{First 10 U\PYZus{}i:}\PY{l+s+s2}{\PYZdq{}}\PY{p}{,} \PY{n}{U}\PY{p}{[}\PY{p}{:}\PY{l+m+mi}{10}\PY{p}{]}\PY{p}{)}
\end{Verbatim}
\end{tcolorbox}

    \begin{Verbatim}[commandchars=\\\{\}]
First 10 X\_i: [1, 5, 25, 125, 625, 1077, 1289, 301, 1505, 1381]
First 10 U\_i: [4.88281250e-04 2.44140625e-03 1.22070312e-02 6.10351562e-02
 3.05175781e-01 5.25878906e-01 6.29394531e-01 1.46972656e-01
 7.34863281e-01 6.74316406e-01]
    \end{Verbatim}

    \subsection*{Part (b)}

Now we plot the time series of \(10^4\) successive random numbers
\(U_i\).

    \begin{tcolorbox}[breakable, size=fbox, boxrule=1pt, pad at break*=1mm,colback=cellbackground, colframe=cellborder]
\prompt{In}{incolor}{10}{\boxspacing}
\begin{Verbatim}[commandchars=\\\{\}]
\PY{n}{i} \PY{o}{=} \PY{n}{np}\PY{o}{.}\PY{n}{linspace}\PY{p}{(}\PY{l+m+mi}{1}\PY{p}{,} \PY{l+m+mi}{10}\PY{o}{*}\PY{o}{*}\PY{l+m+mi}{4}\PY{o}{+}\PY{l+m+mi}{1}\PY{p}{,} \PY{l+m+mi}{10}\PY{o}{*}\PY{o}{*}\PY{l+m+mi}{4}\PY{p}{)}

\PY{n}{plt}\PY{o}{.}\PY{n}{plot}\PY{p}{(}\PY{n}{i}\PY{p}{,} \PY{n}{U}\PY{p}{[}\PY{l+m+mi}{1}\PY{p}{:}\PY{l+m+mi}{10}\PY{o}{*}\PY{o}{*}\PY{l+m+mi}{4}\PY{o}{+}\PY{l+m+mi}{1}\PY{p}{]}\PY{p}{,} \PY{l+s+s2}{\PYZdq{}}\PY{l+s+s2}{.}\PY{l+s+s2}{\PYZdq{}}\PY{p}{,} \PY{n}{markersize}\PY{o}{=}\PY{l+m+mi}{2}\PY{p}{)}
\PY{n}{plt}\PY{o}{.}\PY{n}{xlabel}\PY{p}{(}\PY{l+s+s2}{\PYZdq{}}\PY{l+s+s2}{i}\PY{l+s+s2}{\PYZdq{}}\PY{p}{)}
\PY{n}{plt}\PY{o}{.}\PY{n}{ylabel}\PY{p}{(}\PY{l+s+sa}{r}\PY{l+s+s2}{\PYZdq{}}\PY{l+s+s2}{\PYZdl{}U\PYZus{}i\PYZdl{}}\PY{l+s+s2}{\PYZdq{}}\PY{p}{)}
\PY{n}{plt}\PY{o}{.}\PY{n}{show}\PY{p}{(}\PY{p}{)}
\end{Verbatim}
\end{tcolorbox}

    \begin{center}
    \adjustimage{max size={0.9\linewidth}{0.9\paperheight}}{CompSim_HW4_files/CompSim_HW4_5_0.png}
    \end{center}
    { \hspace*{\fill} \\}
    
    Figure 1. Time series of \(10^4\) successive random numbers \(U_i\) of
the LCG

    The time series does not look random, we see a regular pattern that
repeats itself. The period \(\tau\) is the number of steps after which
the generator returns to its starting value \(X_0\). Since \(X_{i+1}\)
depends only on \(X_i\), from then on the whole sequence repeats itself.
We determine \(\tau\) by running through the sequence until \(X_0\)
appears again.

    \begin{tcolorbox}[breakable, size=fbox, boxrule=1pt, pad at break*=1mm,colback=cellbackground, colframe=cellborder]
\prompt{In}{incolor}{11}{\boxspacing}
\begin{Verbatim}[commandchars=\\\{\}]
\PY{n}{tau} \PY{o}{=} \PY{l+m+mi}{1}
\PY{k}{while} \PY{n}{X}\PY{p}{[}\PY{n}{tau}\PY{p}{]} \PY{o}{!=} \PY{n}{X\PYZus{}0}\PY{p}{:}
    \PY{n}{tau} \PY{o}{+}\PY{o}{=} \PY{l+m+mi}{1}

\PY{n+nb}{print}\PY{p}{(}\PY{l+s+s2}{\PYZdq{}}\PY{l+s+s2}{Period tau =}\PY{l+s+s2}{\PYZdq{}}\PY{p}{,} \PY{n}{tau}\PY{p}{)}
\end{Verbatim}
\end{tcolorbox}

    \begin{Verbatim}[commandchars=\\\{\}]
Period tau = 512
    \end{Verbatim}

    To see the period better, we plot only the first three periods and mark
\(i=\tau, 2\tau, 3\tau\).

    \begin{tcolorbox}[breakable, size=fbox, boxrule=1pt, pad at break*=1mm,colback=cellbackground, colframe=cellborder]
\prompt{In}{incolor}{12}{\boxspacing}
\begin{Verbatim}[commandchars=\\\{\}]
\PY{n}{i} \PY{o}{=} \PY{n}{np}\PY{o}{.}\PY{n}{linspace}\PY{p}{(}\PY{l+m+mi}{1}\PY{p}{,} \PY{l+m+mi}{3}\PY{o}{*}\PY{n}{tau}\PY{o}{+}\PY{l+m+mi}{1}\PY{p}{,} \PY{l+m+mi}{3}\PY{o}{*}\PY{n}{tau}\PY{p}{)}

\PY{n}{plt}\PY{o}{.}\PY{n}{plot}\PY{p}{(}\PY{n}{i}\PY{p}{,} \PY{n}{U}\PY{p}{[}\PY{l+m+mi}{1}\PY{p}{:}\PY{l+m+mi}{3}\PY{o}{*}\PY{n}{tau}\PY{o}{+}\PY{l+m+mi}{1}\PY{p}{]}\PY{p}{,} \PY{l+s+s2}{\PYZdq{}}\PY{l+s+s2}{.}\PY{l+s+s2}{\PYZdq{}}\PY{p}{,} \PY{n}{markersize}\PY{o}{=}\PY{l+m+mi}{3}\PY{p}{)}
\PY{k}{for} \PY{n}{k} \PY{o+ow}{in} \PY{p}{[}\PY{l+m+mi}{1}\PY{p}{,} \PY{l+m+mi}{2}\PY{p}{,} \PY{l+m+mi}{3}\PY{p}{]}\PY{p}{:}
    \PY{n}{plt}\PY{o}{.}\PY{n}{axvline}\PY{p}{(}\PY{n}{k}\PY{o}{*}\PY{n}{tau}\PY{p}{,} \PY{n}{linestyle}\PY{o}{=}\PY{l+s+s2}{\PYZdq{}}\PY{l+s+s2}{\PYZhy{}\PYZhy{}}\PY{l+s+s2}{\PYZdq{}}\PY{p}{,} \PY{n}{color}\PY{o}{=}\PY{l+s+s2}{\PYZdq{}}\PY{l+s+s2}{red}\PY{l+s+s2}{\PYZdq{}}\PY{p}{)}
\PY{n}{plt}\PY{o}{.}\PY{n}{xlabel}\PY{p}{(}\PY{l+s+s2}{\PYZdq{}}\PY{l+s+s2}{i}\PY{l+s+s2}{\PYZdq{}}\PY{p}{)}
\PY{n}{plt}\PY{o}{.}\PY{n}{ylabel}\PY{p}{(}\PY{l+s+sa}{r}\PY{l+s+s2}{\PYZdq{}}\PY{l+s+s2}{\PYZdl{}U\PYZus{}i\PYZdl{}}\PY{l+s+s2}{\PYZdq{}}\PY{p}{)}
\PY{n}{plt}\PY{o}{.}\PY{n}{show}\PY{p}{(}\PY{p}{)}
\end{Verbatim}
\end{tcolorbox}

    \begin{center}
    \adjustimage{max size={0.9\linewidth}{0.9\paperheight}}{CompSim_HW4_files/CompSim_HW4_10_0.png}
    \end{center}
    { \hspace*{\fill} \\}
    
    Figure 2. First three periods of the time series. The red dashed lines
are at \(i=\tau, 2\tau, 3\tau\)

    So the period is \(\tau=512\).

The period tells us how many different random numbers the generator can
produce before it repeats itself.

    \subsection*{Part (c)}

Now we compute the average and the estimator for the variance

\[
\bar{X}=\frac{1}{N}\sum_{i=1}^{N}U_i, \qquad \hat{\sigma}_X^2=\overline{X^2}-\bar{X}^2 \tag{3}
\]

for \(N=\tau, 10\tau, 100\tau\). For equally distributed numbers in
\([0,1)\) the exact values are \(1/2\) and \(1/12\).

    \begin{tcolorbox}[breakable, size=fbox, boxrule=1pt, pad at break*=1mm,colback=cellbackground, colframe=cellborder]
\prompt{In}{incolor}{13}{\boxspacing}
\begin{Verbatim}[commandchars=\\\{\}]
\PY{n+nb}{print}\PY{p}{(}\PY{l+s+s2}{\PYZdq{}}\PY{l+s+s2}{EXACT VALUES: mean =}\PY{l+s+s2}{\PYZdq{}}\PY{p}{,} \PY{l+m+mi}{1}\PY{o}{/}\PY{l+m+mi}{2}\PY{p}{,} \PY{l+s+s2}{\PYZdq{}}\PY{l+s+s2}{, variance =}\PY{l+s+s2}{\PYZdq{}}\PY{p}{,} \PY{l+m+mi}{1}\PY{o}{/}\PY{l+m+mi}{12}\PY{p}{)}

\PY{k}{for} \PY{n}{N} \PY{o+ow}{in} \PY{p}{[}\PY{n}{tau}\PY{p}{,} \PY{l+m+mi}{10}\PY{o}{*}\PY{n}{tau}\PY{p}{,} \PY{l+m+mi}{100}\PY{o}{*}\PY{n}{tau}\PY{p}{]}\PY{p}{:}
    \PY{n}{X\PYZus{}bar} \PY{o}{=} \PY{n}{np}\PY{o}{.}\PY{n}{mean}\PY{p}{(}\PY{n}{U}\PY{p}{[}\PY{l+m+mi}{1}\PY{p}{:}\PY{n}{N}\PY{o}{+}\PY{l+m+mi}{1}\PY{p}{]}\PY{p}{)}
    \PY{n}{X2\PYZus{}bar} \PY{o}{=} \PY{n}{np}\PY{o}{.}\PY{n}{mean}\PY{p}{(}\PY{n}{U}\PY{p}{[}\PY{l+m+mi}{1}\PY{p}{:}\PY{n}{N}\PY{o}{+}\PY{l+m+mi}{1}\PY{p}{]}\PY{o}{*}\PY{o}{*}\PY{l+m+mi}{2}\PY{p}{)}
    \PY{n+nb}{print}\PY{p}{(}\PY{l+s+sa}{f}\PY{l+s+s2}{\PYZdq{}}\PY{l+s+s2}{For N=}\PY{l+s+si}{\PYZob{}}\PY{n}{N}\PY{l+s+si}{\PYZcb{}}\PY{l+s+s2}{: mean = }\PY{l+s+si}{\PYZob{}}\PY{n}{X\PYZus{}bar}\PY{l+s+si}{\PYZcb{}}\PY{l+s+s2}{, variance = }\PY{l+s+si}{\PYZob{}}\PY{n}{X2\PYZus{}bar}\PY{+w}{ }\PY{o}{\PYZhy{}}\PY{+w}{ }\PY{n}{X\PYZus{}bar}\PY{o}{*}\PY{o}{*}\PY{l+m+mi}{2}\PY{l+s+si}{\PYZcb{}}\PY{l+s+s2}{\PYZdq{}}\PY{p}{)}
\end{Verbatim}
\end{tcolorbox}

    \begin{Verbatim}[commandchars=\\\{\}]
EXACT VALUES: mean = 0.5 , variance = 0.08333333333333333
For N=512: mean = 0.49951171875, variance = 0.08333301544189453
For N=5120: mean = 0.49951171875, variance = 0.08333301544189453
For N=51200: mean = 0.49951171875, variance = 0.08333301544189453
    \end{Verbatim}

    The results are the same for all three values of \(N\) (apart from
rounding errors in the last digits). For independent random numbers the
error of the average decreases with \(1/\sqrt{N}\). Here, however, after
\(\tau\) numbers the same numbers come again. For \(N=10\tau\) and
\(N=100\tau\) we just average \(10\) and \(100\) times over the same
\(512\) numbers, so we get exactly the same result as for \(N=\tau\) and
no new information.

The difference to \(1/2\) and \(1/12\) is not a statistical error but
comes from the generator itself. One period contains the \(512\) numbers
\((1+4k)/2048\), \(k=0,\dots,511\), and their average is

\[
\bar{X}_\tau=\frac{1023}{2048}=0.49951171875 \tag{4}
\]

which is exactly the value we found. This can only be improved by a
generator with a much larger period.

    \subsection*{Part (d)}

The running average is

\[
\bar{X}(n)=\frac{1}{n}\sum_{i=1}^{n}U_i \tag{5}
\]

and we plot it for \(n\in[30000,40000]\).

    \begin{tcolorbox}[breakable, size=fbox, boxrule=1pt, pad at break*=1mm,colback=cellbackground, colframe=cellborder]
\prompt{In}{incolor}{14}{\boxspacing}
\begin{Verbatim}[commandchars=\\\{\}]
\PY{n}{running\PYZus{}average} \PY{o}{=} \PY{n}{np}\PY{o}{.}\PY{n}{cumsum}\PY{p}{(}\PY{n}{U}\PY{p}{[}\PY{l+m+mi}{1}\PY{p}{:}\PY{p}{]}\PY{p}{)}\PY{o}{/}\PY{n}{np}\PY{o}{.}\PY{n}{arange}\PY{p}{(}\PY{l+m+mi}{1}\PY{p}{,} \PY{n+nb}{len}\PY{p}{(}\PY{n}{U}\PY{p}{)}\PY{p}{)}

\PY{n}{n} \PY{o}{=} \PY{n}{np}\PY{o}{.}\PY{n}{arange}\PY{p}{(}\PY{l+m+mi}{30000}\PY{p}{,} \PY{l+m+mi}{40001}\PY{p}{)}

\PY{n}{plt}\PY{o}{.}\PY{n}{plot}\PY{p}{(}\PY{n}{n}\PY{p}{,} \PY{n}{running\PYZus{}average}\PY{p}{[}\PY{n}{n}\PY{o}{\PYZhy{}}\PY{l+m+mi}{1}\PY{p}{]}\PY{p}{,} \PY{n}{label}\PY{o}{=}\PY{l+s+sa}{r}\PY{l+s+s2}{\PYZdq{}}\PY{l+s+s2}{\PYZdl{}}\PY{l+s+s2}{\PYZbs{}}\PY{l+s+s2}{bar}\PY{l+s+si}{\PYZob{}X\PYZcb{}}\PY{l+s+s2}{(n)\PYZdl{}}\PY{l+s+s2}{\PYZdq{}}\PY{p}{)}
\PY{n}{plt}\PY{o}{.}\PY{n}{axhline}\PY{p}{(}\PY{l+m+mi}{1}\PY{o}{/}\PY{l+m+mi}{2}\PY{p}{,} \PY{n}{linestyle}\PY{o}{=}\PY{l+s+s2}{\PYZdq{}}\PY{l+s+s2}{\PYZhy{}\PYZhy{}}\PY{l+s+s2}{\PYZdq{}}\PY{p}{,} \PY{n}{color}\PY{o}{=}\PY{l+s+s2}{\PYZdq{}}\PY{l+s+s2}{red}\PY{l+s+s2}{\PYZdq{}}\PY{p}{,} \PY{n}{label}\PY{o}{=}\PY{l+s+s2}{\PYZdq{}}\PY{l+s+s2}{1/2}\PY{l+s+s2}{\PYZdq{}}\PY{p}{)}
\PY{n}{plt}\PY{o}{.}\PY{n}{axhline}\PY{p}{(}\PY{l+m+mi}{1023}\PY{o}{/}\PY{l+m+mi}{2048}\PY{p}{,} \PY{n}{linestyle}\PY{o}{=}\PY{l+s+s2}{\PYZdq{}}\PY{l+s+s2}{\PYZhy{}\PYZhy{}}\PY{l+s+s2}{\PYZdq{}}\PY{p}{,} \PY{n}{color}\PY{o}{=}\PY{l+s+s2}{\PYZdq{}}\PY{l+s+s2}{green}\PY{l+s+s2}{\PYZdq{}}\PY{p}{,} \PY{n}{label}\PY{o}{=}\PY{l+s+s2}{\PYZdq{}}\PY{l+s+s2}{1023/2048}\PY{l+s+s2}{\PYZdq{}}\PY{p}{)}
\PY{n}{plt}\PY{o}{.}\PY{n}{xlabel}\PY{p}{(}\PY{l+s+s2}{\PYZdq{}}\PY{l+s+s2}{n}\PY{l+s+s2}{\PYZdq{}}\PY{p}{)}
\PY{n}{plt}\PY{o}{.}\PY{n}{ylabel}\PY{p}{(}\PY{l+s+sa}{r}\PY{l+s+s2}{\PYZdq{}}\PY{l+s+s2}{\PYZdl{}}\PY{l+s+s2}{\PYZbs{}}\PY{l+s+s2}{bar}\PY{l+s+si}{\PYZob{}X\PYZcb{}}\PY{l+s+s2}{(n)\PYZdl{}}\PY{l+s+s2}{\PYZdq{}}\PY{p}{)}
\PY{n}{plt}\PY{o}{.}\PY{n}{legend}\PY{p}{(}\PY{p}{)}
\PY{n}{plt}\PY{o}{.}\PY{n}{show}\PY{p}{(}\PY{p}{)}
\end{Verbatim}
\end{tcolorbox}

    \begin{center}
    \adjustimage{max size={0.9\linewidth}{0.9\paperheight}}{CompSim_HW4_files/CompSim_HW4_17_0.png}
    \end{center}
    { \hspace*{\fill} \\}
    
    Figure 3. Running average \(\bar{X}(n)\) for \(n\in[30000,40000]\)

    \subsection*{Part (e)}

Finally we plot \(x=U_i\) against \(y=U_{i+1}\) for \(i=1,3,5,\dots\)

    \begin{tcolorbox}[breakable, size=fbox, boxrule=1pt, pad at break*=1mm,colback=cellbackground, colframe=cellborder]
\prompt{In}{incolor}{15}{\boxspacing}
\begin{Verbatim}[commandchars=\\\{\}]
\PY{n}{x} \PY{o}{=} \PY{n}{U}\PY{p}{[}\PY{l+m+mi}{1}\PY{p}{:}\PY{l+m+mi}{10}\PY{o}{*}\PY{o}{*}\PY{l+m+mi}{4}\PY{p}{:}\PY{l+m+mi}{2}\PY{p}{]}
\PY{n}{y} \PY{o}{=} \PY{n}{U}\PY{p}{[}\PY{l+m+mi}{2}\PY{p}{:}\PY{l+m+mi}{10}\PY{o}{*}\PY{o}{*}\PY{l+m+mi}{4}\PY{o}{+}\PY{l+m+mi}{1}\PY{p}{:}\PY{l+m+mi}{2}\PY{p}{]}

\PY{n}{plt}\PY{o}{.}\PY{n}{plot}\PY{p}{(}\PY{n}{x}\PY{p}{,} \PY{n}{y}\PY{p}{,} \PY{l+s+s2}{\PYZdq{}}\PY{l+s+s2}{.}\PY{l+s+s2}{\PYZdq{}}\PY{p}{,} \PY{n}{markersize}\PY{o}{=}\PY{l+m+mi}{4}\PY{p}{)}
\PY{n}{plt}\PY{o}{.}\PY{n}{xlabel}\PY{p}{(}\PY{l+s+sa}{r}\PY{l+s+s2}{\PYZdq{}}\PY{l+s+s2}{\PYZdl{}x=U\PYZus{}i\PYZdl{}}\PY{l+s+s2}{\PYZdq{}}\PY{p}{)}
\PY{n}{plt}\PY{o}{.}\PY{n}{ylabel}\PY{p}{(}\PY{l+s+sa}{r}\PY{l+s+s2}{\PYZdq{}}\PY{l+s+s2}{\PYZdl{}y=U\PYZus{}}\PY{l+s+s2}{\PYZob{}}\PY{l+s+s2}{i+1\PYZcb{}\PYZdl{}}\PY{l+s+s2}{\PYZdq{}}\PY{p}{)}
\PY{n}{plt}\PY{o}{.}\PY{n}{show}\PY{p}{(}\PY{p}{)}

\PY{n+nb}{print}\PY{p}{(}\PY{l+s+s2}{\PYZdq{}}\PY{l+s+s2}{Number of points:}\PY{l+s+s2}{\PYZdq{}}\PY{p}{,} \PY{n+nb}{len}\PY{p}{(}\PY{n}{x}\PY{p}{)}\PY{p}{)}
\PY{n+nb}{print}\PY{p}{(}\PY{l+s+s2}{\PYZdq{}}\PY{l+s+s2}{Number of different points:}\PY{l+s+s2}{\PYZdq{}}\PY{p}{,} \PY{n+nb}{len}\PY{p}{(}\PY{n+nb}{set}\PY{p}{(}\PY{n+nb}{zip}\PY{p}{(}\PY{n}{x}\PY{p}{,} \PY{n}{y}\PY{p}{)}\PY{p}{)}\PY{p}{)}\PY{p}{)}
\end{Verbatim}
\end{tcolorbox}

    \begin{center}
    \adjustimage{max size={0.9\linewidth}{0.9\paperheight}}{CompSim_HW4_files/CompSim_HW4_20_0.png}
    \end{center}
    { \hspace*{\fill} \\}
    
    \begin{Verbatim}[commandchars=\\\{\}]
Number of points: 5000
Number of different points: 256
    \end{Verbatim}

    Figure 4. \(U_{i+1}\) against \(U_i\) for \(i=1,3,5,\dots,9999\)

    \section*{Problem 8: Statement}\label{problem-8-statement}

Write a Monte Carlo computer program for the simulation of the 1D Ising
model with \(N = 20, 100, 1000, 10000\) spins and free boundary
conditions. Use for the spin updates the Metropolis algorithm.

Determine the mean energy, specific heat, magnetisation and
susceptibility for different temperatures in the inverse temperature
range \(\beta=J/(k_bT) =0,...,2.\)

Check your results by comparison with the analytical results for the
mean energy and specific heat of Problem 1 respectively for N=20 spins
and \emph{all} quantities by means of the exact enumerations of Problems
3 and 5.

Measure also the acceptance rate of the spin-flip proposals.

    \section*{Problem 8: Solution}\label{problem-8-solution}

    We simulate the 1D Ising model with free boundary conditions and
\(h=0\). As in Homework 1, the energy is

\[
E=-J\sum_{i=1}^{N-1}s_is_{i+1} \tag{7}
\]

and we set \(J=k_B=1\), so that \(\beta=1/T\).

In the Metropolis algorithm we choose a spin \(i\) at random and
calculate the energy change if this spin is flipped

\[
\Delta E=2Js_i(s_{i-1}+s_{i+1}) \tag{8}
\]

where for the first and the last spin the missing neighbour is left out
(free boundary conditions). If \(\Delta E\le0\) the flip is accepted,
otherwise it is accepted with probability \(e^{-\beta\Delta E}\). \(N\)
such steps are one sweep. We start from a random configuration, do
\(10^4\) sweeps without measuring to reach equilibrium, and then measure
\(E\) and \(M=\sum_i s_i\) after every sweep.

From the measurements we calculate the per spin quantities

\[
e=\frac{\langle E\rangle}{N}, \qquad c=\frac{\beta^2}{N}\left(\langle E^2\rangle-\langle E\rangle^2\right) \tag{9}
\]

\[
m=\frac{\langle |M|\rangle}{N}, \qquad \chi=\frac{\beta}{N}\left(\langle M^2\rangle-\langle M\rangle^2\right)=\frac{\beta}{N}\langle M^2\rangle \tag{10}
\]

For \(h=0\) we have \(\langle M\rangle=0\) by symmetry, therefore we use
\(\langle |M|\rangle\) for the magnetisation. The acceptance rate is the
number of accepted flips divided by the number of proposed flips.

    \begin{tcolorbox}[breakable, size=fbox, boxrule=1pt, pad at break*=1mm,colback=cellbackground, colframe=cellborder]
\prompt{In}{incolor}{16}{\boxspacing}
\begin{Verbatim}[commandchars=\\\{\}]
\PY{k+kn}{import} \PY{n+nn}{numpy} \PY{k}{as} \PY{n+nn}{np}
\PY{k+kn}{import} \PY{n+nn}{math}
\PY{k+kn}{import} \PY{n+nn}{time}
\PY{k+kn}{from} \PY{n+nn}{numba} \PY{k+kn}{import} \PY{n}{njit}\PY{p}{,} \PY{n}{prange}
\PY{k+kn}{from} \PY{n+nn}{matplotlib} \PY{k+kn}{import} \PY{n}{pyplot} \PY{k}{as} \PY{n}{plt}

\PY{n}{J} \PY{o}{=} \PY{l+m+mf}{1.0}
\PY{n}{k\PYZus{}b} \PY{o}{=} \PY{l+m+mf}{1.0}

\PY{n+nd}{@njit}
\PY{k}{def} \PY{n+nf}{metropolis}\PY{p}{(}\PY{n}{N}\PY{p}{,} \PY{n}{beta}\PY{p}{,} \PY{n}{n\PYZus{}therm}\PY{p}{,} \PY{n}{n\PYZus{}sweeps}\PY{p}{,} \PY{n}{seed}\PY{p}{)}\PY{p}{:}
    \PY{n}{np}\PY{o}{.}\PY{n}{random}\PY{o}{.}\PY{n}{seed}\PY{p}{(}\PY{n}{seed}\PY{p}{)}

    \PY{c+c1}{\PYZsh{} random start configuration}
    \PY{n}{s} \PY{o}{=} \PY{n}{np}\PY{o}{.}\PY{n}{ones}\PY{p}{(}\PY{n}{N}\PY{p}{)}
    \PY{k}{for} \PY{n}{i} \PY{o+ow}{in} \PY{n+nb}{range}\PY{p}{(}\PY{n}{N}\PY{p}{)}\PY{p}{:}
        \PY{k}{if} \PY{n}{np}\PY{o}{.}\PY{n}{random}\PY{o}{.}\PY{n}{random}\PY{p}{(}\PY{p}{)} \PY{o}{\PYZlt{}} \PY{l+m+mf}{0.5}\PY{p}{:}
            \PY{n}{s}\PY{p}{[}\PY{n}{i}\PY{p}{]} \PY{o}{=} \PY{o}{\PYZhy{}}\PY{l+m+mi}{1}

    \PY{n}{E} \PY{o}{=} \PY{l+m+mf}{0.0}
    \PY{k}{for} \PY{n}{i} \PY{o+ow}{in} \PY{n+nb}{range}\PY{p}{(}\PY{n}{N}\PY{o}{\PYZhy{}}\PY{l+m+mi}{1}\PY{p}{)}\PY{p}{:}
        \PY{n}{E} \PY{o}{+}\PY{o}{=} \PY{o}{\PYZhy{}}\PY{n}{J}\PY{o}{*}\PY{n}{s}\PY{p}{[}\PY{n}{i}\PY{p}{]}\PY{o}{*}\PY{n}{s}\PY{p}{[}\PY{n}{i}\PY{o}{+}\PY{l+m+mi}{1}\PY{p}{]}
    \PY{n}{M} \PY{o}{=} \PY{n}{np}\PY{o}{.}\PY{n}{sum}\PY{p}{(}\PY{n}{s}\PY{p}{)}

    \PY{n}{E\PYZus{}sum} \PY{o}{=} \PY{l+m+mf}{0.0}
    \PY{n}{E2\PYZus{}sum} \PY{o}{=} \PY{l+m+mf}{0.0}
    \PY{n}{M\PYZus{}sum} \PY{o}{=} \PY{l+m+mf}{0.0}
    \PY{n}{M2\PYZus{}sum} \PY{o}{=} \PY{l+m+mf}{0.0}
    \PY{n}{accepted} \PY{o}{=} \PY{l+m+mi}{0}

    \PY{k}{for} \PY{n}{sweep} \PY{o+ow}{in} \PY{n+nb}{range}\PY{p}{(}\PY{n}{n\PYZus{}therm} \PY{o}{+} \PY{n}{n\PYZus{}sweeps}\PY{p}{)}\PY{p}{:}
        \PY{k}{for} \PY{n}{step} \PY{o+ow}{in} \PY{n+nb}{range}\PY{p}{(}\PY{n}{N}\PY{p}{)}\PY{p}{:}
            \PY{n}{i} \PY{o}{=} \PY{n}{np}\PY{o}{.}\PY{n}{random}\PY{o}{.}\PY{n}{randint}\PY{p}{(}\PY{n}{N}\PY{p}{)}
            \PY{n}{neighbours} \PY{o}{=} \PY{l+m+mf}{0.0}
            \PY{k}{if} \PY{n}{i} \PY{o}{\PYZgt{}} \PY{l+m+mi}{0}\PY{p}{:}
                \PY{n}{neighbours} \PY{o}{+}\PY{o}{=} \PY{n}{s}\PY{p}{[}\PY{n}{i}\PY{o}{\PYZhy{}}\PY{l+m+mi}{1}\PY{p}{]}
            \PY{k}{if} \PY{n}{i} \PY{o}{\PYZlt{}} \PY{n}{N}\PY{o}{\PYZhy{}}\PY{l+m+mi}{1}\PY{p}{:}
                \PY{n}{neighbours} \PY{o}{+}\PY{o}{=} \PY{n}{s}\PY{p}{[}\PY{n}{i}\PY{o}{+}\PY{l+m+mi}{1}\PY{p}{]}
            \PY{n}{dE} \PY{o}{=} \PY{l+m+mi}{2}\PY{o}{*}\PY{n}{J}\PY{o}{*}\PY{n}{s}\PY{p}{[}\PY{n}{i}\PY{p}{]}\PY{o}{*}\PY{n}{neighbours}
            \PY{k}{if} \PY{n}{dE} \PY{o}{\PYZlt{}}\PY{o}{=} \PY{l+m+mi}{0} \PY{o+ow}{or} \PY{n}{np}\PY{o}{.}\PY{n}{random}\PY{o}{.}\PY{n}{random}\PY{p}{(}\PY{p}{)} \PY{o}{\PYZlt{}} \PY{n}{np}\PY{o}{.}\PY{n}{exp}\PY{p}{(}\PY{o}{\PYZhy{}}\PY{n}{beta}\PY{o}{*}\PY{n}{dE}\PY{p}{)}\PY{p}{:}
                \PY{n}{s}\PY{p}{[}\PY{n}{i}\PY{p}{]} \PY{o}{=} \PY{o}{\PYZhy{}}\PY{n}{s}\PY{p}{[}\PY{n}{i}\PY{p}{]}
                \PY{n}{E} \PY{o}{+}\PY{o}{=} \PY{n}{dE}
                \PY{n}{M} \PY{o}{+}\PY{o}{=} \PY{l+m+mi}{2}\PY{o}{*}\PY{n}{s}\PY{p}{[}\PY{n}{i}\PY{p}{]}
                \PY{k}{if} \PY{n}{sweep} \PY{o}{\PYZgt{}}\PY{o}{=} \PY{n}{n\PYZus{}therm}\PY{p}{:}
                    \PY{n}{accepted} \PY{o}{+}\PY{o}{=} \PY{l+m+mi}{1}
        \PY{k}{if} \PY{n}{sweep} \PY{o}{\PYZgt{}}\PY{o}{=} \PY{n}{n\PYZus{}therm}\PY{p}{:}
            \PY{n}{E\PYZus{}sum} \PY{o}{+}\PY{o}{=} \PY{n}{E}
            \PY{n}{E2\PYZus{}sum} \PY{o}{+}\PY{o}{=} \PY{n}{E}\PY{o}{*}\PY{n}{E}
            \PY{n}{M\PYZus{}sum} \PY{o}{+}\PY{o}{=} \PY{n+nb}{abs}\PY{p}{(}\PY{n}{M}\PY{p}{)}
            \PY{n}{M2\PYZus{}sum} \PY{o}{+}\PY{o}{=} \PY{n}{M}\PY{o}{*}\PY{n}{M}

    \PY{k}{return} \PY{n}{E\PYZus{}sum}\PY{o}{/}\PY{n}{n\PYZus{}sweeps}\PY{p}{,} \PY{n}{E2\PYZus{}sum}\PY{o}{/}\PY{n}{n\PYZus{}sweeps}\PY{p}{,} \PY{n}{M\PYZus{}sum}\PY{o}{/}\PY{n}{n\PYZus{}sweeps}\PY{p}{,} \PY{n}{M2\PYZus{}sum}\PY{o}{/}\PY{n}{n\PYZus{}sweeps}\PY{p}{,} \PY{n}{accepted}\PY{o}{/}\PY{p}{(}\PY{n}{N}\PY{o}{*}\PY{n}{n\PYZus{}sweeps}\PY{p}{)}

\PY{n+nd}{@njit}\PY{p}{(}\PY{n}{parallel}\PY{o}{=}\PY{k+kc}{True}\PY{p}{)}
\PY{k}{def} \PY{n+nf}{simulate}\PY{p}{(}\PY{n}{N}\PY{p}{,} \PY{n}{betas}\PY{p}{,} \PY{n}{n\PYZus{}therm}\PY{p}{,} \PY{n}{n\PYZus{}sweeps}\PY{p}{)}\PY{p}{:}
    \PY{n}{results} \PY{o}{=} \PY{n}{np}\PY{o}{.}\PY{n}{zeros}\PY{p}{(}\PY{p}{(}\PY{n+nb}{len}\PY{p}{(}\PY{n}{betas}\PY{p}{)}\PY{p}{,} \PY{l+m+mi}{5}\PY{p}{)}\PY{p}{)}
    \PY{k}{for} \PY{n}{k} \PY{o+ow}{in} \PY{n}{prange}\PY{p}{(}\PY{n+nb}{len}\PY{p}{(}\PY{n}{betas}\PY{p}{)}\PY{p}{)}\PY{p}{:}
        \PY{n}{E}\PY{p}{,} \PY{n}{E2}\PY{p}{,} \PY{n}{M\PYZus{}abs}\PY{p}{,} \PY{n}{M2}\PY{p}{,} \PY{n}{acc} \PY{o}{=} \PY{n}{metropolis}\PY{p}{(}\PY{n}{N}\PY{p}{,} \PY{n}{betas}\PY{p}{[}\PY{n}{k}\PY{p}{]}\PY{p}{,} \PY{n}{n\PYZus{}therm}\PY{p}{,} \PY{n}{n\PYZus{}sweeps}\PY{p}{,} \PY{l+m+mi}{1000}\PY{o}{*}\PY{n}{N} \PY{o}{+} \PY{n}{k}\PY{p}{)}
        \PY{n}{results}\PY{p}{[}\PY{n}{k}\PY{p}{,} \PY{l+m+mi}{0}\PY{p}{]} \PY{o}{=} \PY{n}{E}
        \PY{n}{results}\PY{p}{[}\PY{n}{k}\PY{p}{,} \PY{l+m+mi}{1}\PY{p}{]} \PY{o}{=} \PY{n}{E2}
        \PY{n}{results}\PY{p}{[}\PY{n}{k}\PY{p}{,} \PY{l+m+mi}{2}\PY{p}{]} \PY{o}{=} \PY{n}{M\PYZus{}abs}
        \PY{n}{results}\PY{p}{[}\PY{n}{k}\PY{p}{,} \PY{l+m+mi}{3}\PY{p}{]} \PY{o}{=} \PY{n}{M2}
        \PY{n}{results}\PY{p}{[}\PY{n}{k}\PY{p}{,} \PY{l+m+mi}{4}\PY{p}{]} \PY{o}{=} \PY{n}{acc}
    \PY{k}{return} \PY{n}{results}
\end{Verbatim}
\end{tcolorbox}

    Now we run the simulation for \(N=20,100,1000,10000\) and
\(\beta=0,0.1,\dots,2\).

    \begin{tcolorbox}[breakable, size=fbox, boxrule=1pt, pad at break*=1mm,colback=cellbackground, colframe=cellborder]
\prompt{In}{incolor}{17}{\boxspacing}
\begin{Verbatim}[commandchars=\\\{\}]
\PY{n}{betas} \PY{o}{=} \PY{n}{np}\PY{o}{.}\PY{n}{linspace}\PY{p}{(}\PY{l+m+mi}{0}\PY{p}{,} \PY{l+m+mi}{2}\PY{p}{,} \PY{l+m+mi}{21}\PY{p}{)}
\PY{n}{N\PYZus{}values} \PY{o}{=} \PY{p}{[}\PY{l+m+mi}{20}\PY{p}{,} \PY{l+m+mi}{100}\PY{p}{,} \PY{l+m+mi}{1000}\PY{p}{,} \PY{l+m+mi}{10000}\PY{p}{]}

\PY{n}{n\PYZus{}therm} \PY{o}{=} \PY{l+m+mi}{10}\PY{o}{*}\PY{o}{*}\PY{l+m+mi}{4}
\PY{n}{n\PYZus{}sweeps} \PY{o}{=} \PY{p}{\PYZob{}}\PY{l+m+mi}{20}\PY{p}{:} \PY{l+m+mi}{400000}\PY{p}{,} \PY{l+m+mi}{100}\PY{p}{:} \PY{l+m+mi}{400000}\PY{p}{,} \PY{l+m+mi}{1000}\PY{p}{:} \PY{l+m+mi}{200000}\PY{p}{,} \PY{l+m+mi}{10000}\PY{p}{:} \PY{l+m+mi}{100000}\PY{p}{\PYZcb{}}

\PY{n}{e\PYZus{}MC} \PY{o}{=} \PY{p}{\PYZob{}}\PY{p}{\PYZcb{}}
\PY{n}{c\PYZus{}MC} \PY{o}{=} \PY{p}{\PYZob{}}\PY{p}{\PYZcb{}}
\PY{n}{m\PYZus{}MC} \PY{o}{=} \PY{p}{\PYZob{}}\PY{p}{\PYZcb{}}
\PY{n}{chi\PYZus{}MC} \PY{o}{=} \PY{p}{\PYZob{}}\PY{p}{\PYZcb{}}
\PY{n}{acc\PYZus{}MC} \PY{o}{=} \PY{p}{\PYZob{}}\PY{p}{\PYZcb{}}

\PY{n+nb}{print}\PY{p}{(}\PY{l+s+s2}{\PYZdq{}}\PY{l+s+s2}{MONTE CARLO SIMULATION:}\PY{l+s+s2}{\PYZdq{}}\PY{p}{)}
\PY{k}{for} \PY{n}{N} \PY{o+ow}{in} \PY{n}{N\PYZus{}values}\PY{p}{:}
    \PY{n}{t\PYZus{}start} \PY{o}{=} \PY{n}{time}\PY{o}{.}\PY{n}{time}\PY{p}{(}\PY{p}{)}
    \PY{n}{results} \PY{o}{=} \PY{n}{simulate}\PY{p}{(}\PY{n}{N}\PY{p}{,} \PY{n}{betas}\PY{p}{,} \PY{n}{n\PYZus{}therm}\PY{p}{,} \PY{n}{n\PYZus{}sweeps}\PY{p}{[}\PY{n}{N}\PY{p}{]}\PY{p}{)}
    \PY{n}{E}\PY{p}{,} \PY{n}{E2}\PY{p}{,} \PY{n}{M\PYZus{}abs}\PY{p}{,} \PY{n}{M2}\PY{p}{,} \PY{n}{acc} \PY{o}{=} \PY{n}{results}\PY{p}{[}\PY{p}{:}\PY{p}{,} \PY{l+m+mi}{0}\PY{p}{]}\PY{p}{,} \PY{n}{results}\PY{p}{[}\PY{p}{:}\PY{p}{,} \PY{l+m+mi}{1}\PY{p}{]}\PY{p}{,} \PY{n}{results}\PY{p}{[}\PY{p}{:}\PY{p}{,} \PY{l+m+mi}{2}\PY{p}{]}\PY{p}{,} \PY{n}{results}\PY{p}{[}\PY{p}{:}\PY{p}{,} \PY{l+m+mi}{3}\PY{p}{]}\PY{p}{,} \PY{n}{results}\PY{p}{[}\PY{p}{:}\PY{p}{,} \PY{l+m+mi}{4}\PY{p}{]}

    \PY{n}{e\PYZus{}MC}\PY{p}{[}\PY{n}{N}\PY{p}{]} \PY{o}{=} \PY{n}{E}\PY{o}{/}\PY{n}{N}
    \PY{n}{c\PYZus{}MC}\PY{p}{[}\PY{n}{N}\PY{p}{]} \PY{o}{=} \PY{n}{betas}\PY{o}{*}\PY{o}{*}\PY{l+m+mi}{2}\PY{o}{*}\PY{p}{(}\PY{n}{E2} \PY{o}{\PYZhy{}} \PY{n}{E}\PY{o}{*}\PY{o}{*}\PY{l+m+mi}{2}\PY{p}{)}\PY{o}{/}\PY{n}{N}
    \PY{n}{m\PYZus{}MC}\PY{p}{[}\PY{n}{N}\PY{p}{]} \PY{o}{=} \PY{n}{M\PYZus{}abs}\PY{o}{/}\PY{n}{N}
    \PY{n}{chi\PYZus{}MC}\PY{p}{[}\PY{n}{N}\PY{p}{]} \PY{o}{=} \PY{n}{betas}\PY{o}{*}\PY{n}{M2}\PY{o}{/}\PY{n}{N}
    \PY{n}{acc\PYZus{}MC}\PY{p}{[}\PY{n}{N}\PY{p}{]} \PY{o}{=} \PY{n}{acc}
    \PY{n+nb}{print}\PY{p}{(}\PY{l+s+sa}{f}\PY{l+s+s2}{\PYZdq{}}\PY{l+s+s2}{For N=}\PY{l+s+si}{\PYZob{}}\PY{n}{N}\PY{l+s+si}{\PYZcb{}}\PY{l+s+s2}{: done in }\PY{l+s+si}{\PYZob{}}\PY{n}{time}\PY{o}{.}\PY{n}{time}\PY{p}{(}\PY{p}{)}\PY{+w}{ }\PY{o}{\PYZhy{}}\PY{+w}{ }\PY{n}{t\PYZus{}start}\PY{l+s+si}{:}\PY{l+s+s2}{.1f}\PY{l+s+si}{\PYZcb{}}\PY{l+s+s2}{ s}\PY{l+s+s2}{\PYZdq{}}\PY{p}{)}
\end{Verbatim}
\end{tcolorbox}

    \begin{Verbatim}[commandchars=\\\{\}]
MONTE CARLO SIMULATION:
    \end{Verbatim}

    \begin{Verbatim}[commandchars=\\\{\}]
OMP: Info \#276: omp\_set\_nested routine deprecated, please use
omp\_set\_max\_active\_levels instead.
    \end{Verbatim}

    \begin{Verbatim}[commandchars=\\\{\}]
For N=20: done in 2.1 s
For N=100: done in 1.3 s
For N=1000: done in 6.0 s
For N=10000: done in 38.1 s
    \end{Verbatim}

    From Problem 1 (Homework 1) we know the exact mean energy and specific
heat of the free chain. Per spin they are

\[
e=-J\frac{N-1}{N}\tanh(\beta J), \qquad c=(\beta J)^2\frac{N-1}{N}\operatorname{sech}^2(\beta J) \tag{11}
\]

For \(N=20\) we determine the density of states \(\Omega(E,M)\) by exact
enumeration as in Problems 3 and 5 and calculate all quantities from it.

    \begin{tcolorbox}[breakable, size=fbox, boxrule=1pt, pad at break*=1mm,colback=cellbackground, colframe=cellborder]
\prompt{In}{incolor}{18}{\boxspacing}
\begin{Verbatim}[commandchars=\\\{\}]
\PY{k}{def} \PY{n+nf}{e\PYZus{}analytic}\PY{p}{(}\PY{n}{N}\PY{p}{,} \PY{n}{beta}\PY{p}{)}\PY{p}{:}
    \PY{k}{return} \PY{o}{\PYZhy{}}\PY{n}{J}\PY{o}{*}\PY{p}{(}\PY{n}{N}\PY{o}{\PYZhy{}}\PY{l+m+mi}{1}\PY{p}{)}\PY{o}{/}\PY{n}{N}\PY{o}{*}\PY{n}{np}\PY{o}{.}\PY{n}{tanh}\PY{p}{(}\PY{n}{J}\PY{o}{*}\PY{n}{beta}\PY{p}{)}

\PY{k}{def} \PY{n+nf}{c\PYZus{}analytic}\PY{p}{(}\PY{n}{N}\PY{p}{,} \PY{n}{beta}\PY{p}{)}\PY{p}{:}
    \PY{k}{return} \PY{p}{(}\PY{n}{J}\PY{o}{*}\PY{n}{beta}\PY{p}{)}\PY{o}{*}\PY{o}{*}\PY{l+m+mi}{2}\PY{o}{*}\PY{p}{(}\PY{n}{N}\PY{o}{\PYZhy{}}\PY{l+m+mi}{1}\PY{p}{)}\PY{o}{/}\PY{n}{N}\PY{o}{/}\PY{n}{np}\PY{o}{.}\PY{n}{cosh}\PY{p}{(}\PY{n}{J}\PY{o}{*}\PY{n}{beta}\PY{p}{)}\PY{o}{*}\PY{o}{*}\PY{l+m+mi}{2}
\end{Verbatim}
\end{tcolorbox}

    \begin{tcolorbox}[breakable, size=fbox, boxrule=1pt, pad at break*=1mm,colback=cellbackground, colframe=cellborder]
\prompt{In}{incolor}{19}{\boxspacing}
\begin{Verbatim}[commandchars=\\\{\}]
\PY{k}{def} \PY{n+nf}{exp}\PY{p}{(}\PY{n}{x}\PY{p}{)}\PY{p}{:}
 \PY{k}{return} \PY{n}{math}\PY{o}{.}\PY{n}{exp}\PY{p}{(}\PY{n}{x}\PY{p}{)}

\PY{n}{spins}\PY{o}{=}\PY{p}{[}\PY{o}{\PYZhy{}}\PY{l+m+mi}{1}\PY{p}{,}\PY{l+m+mi}{1}\PY{p}{]}

\PY{n}{distr}\PY{o}{=}\PY{p}{\PYZob{}}\PY{p}{\PYZcb{}}
\PY{k}{for} \PY{n}{s1} \PY{o+ow}{in} \PY{n}{spins}\PY{p}{:}
 \PY{k}{for} \PY{n}{s2} \PY{o+ow}{in} \PY{n}{spins}\PY{p}{:}
  \PY{k}{for} \PY{n}{s3} \PY{o+ow}{in} \PY{n}{spins}\PY{p}{:}
   \PY{k}{for} \PY{n}{s4} \PY{o+ow}{in} \PY{n}{spins}\PY{p}{:}
    \PY{k}{for} \PY{n}{s5} \PY{o+ow}{in} \PY{n}{spins}\PY{p}{:}
     \PY{k}{for} \PY{n}{s6} \PY{o+ow}{in} \PY{n}{spins}\PY{p}{:}
      \PY{k}{for} \PY{n}{s7} \PY{o+ow}{in} \PY{n}{spins}\PY{p}{:}
       \PY{k}{for} \PY{n}{s8} \PY{o+ow}{in} \PY{n}{spins}\PY{p}{:}
        \PY{k}{for} \PY{n}{s9} \PY{o+ow}{in} \PY{n}{spins}\PY{p}{:}
         \PY{k}{for} \PY{n}{s10} \PY{o+ow}{in} \PY{n}{spins}\PY{p}{:}
          \PY{k}{for} \PY{n}{s11} \PY{o+ow}{in} \PY{n}{spins}\PY{p}{:}
           \PY{k}{for} \PY{n}{s12} \PY{o+ow}{in} \PY{n}{spins}\PY{p}{:}
            \PY{k}{for} \PY{n}{s13} \PY{o+ow}{in} \PY{n}{spins}\PY{p}{:}
             \PY{k}{for} \PY{n}{s14} \PY{o+ow}{in} \PY{n}{spins}\PY{p}{:}
              \PY{k}{for} \PY{n}{s15} \PY{o+ow}{in} \PY{n}{spins}\PY{p}{:}
               \PY{k}{for} \PY{n}{s16} \PY{o+ow}{in} \PY{n}{spins}\PY{p}{:}
                \PY{k}{for} \PY{n}{s17} \PY{o+ow}{in} \PY{n}{spins}\PY{p}{:}
                 \PY{k}{for} \PY{n}{s18} \PY{o+ow}{in} \PY{n}{spins}\PY{p}{:}
                  \PY{k}{for} \PY{n}{s19} \PY{o+ow}{in} \PY{n}{spins}\PY{p}{:}
                   \PY{k}{for} \PY{n}{s20} \PY{o+ow}{in} \PY{n}{spins}\PY{p}{:}
                    \PY{n}{E\PYZus{}20}\PY{o}{=}\PY{o}{\PYZhy{}}\PY{n}{J}\PY{o}{*}\PY{p}{(}\PY{n}{s1}\PY{o}{*}\PY{n}{s2}\PY{o}{+}\discretionary{}{\Wrappedafterbreak}{}\PY{n}{s2}\PY{o}{*}\PY{n}{s3}\PY{o}{+}\discretionary{}{\Wrappedafterbreak}{}\PY{n}{s3}\PY{o}{*}\PY{n}{s4}\PY{o}{+}\discretionary{}{\Wrappedafterbreak}{}\PY{n}{s4}\PY{o}{*}\PY{n}{s5}\PY{o}{+}\discretionary{}{\Wrappedafterbreak}{}\PY{n}{s5}\PY{o}{*}\PY{n}{s6}\PY{o}{+}\discretionary{}{\Wrappedafterbreak}{}\PY{n}{s6}\PY{o}{*}\PY{n}{s7}\PY{o}{+}\discretionary{}{\Wrappedafterbreak}{}\PY{n}{s7}\PY{o}{*}\PY{n}{s8}\PY{o}{+}\discretionary{}{\Wrappedafterbreak}{}\PY{n}{s8}\PY{o}{*}\PY{n}{s9}\PY{o}{+}\discretionary{}{\Wrappedafterbreak}{}\PY{n}{s9}\PY{o}{*}\PY{n}{s10}\PY{o}{+}\discretionary{}{\Wrappedafterbreak}{}\PY{n}{s10}\PY{o}{*}\PY{n}{s11}\PY{o}{+}\discretionary{}{\Wrappedafterbreak}{}\PY{n}{s11}\PY{o}{*}\PY{n}{s12}\PY{o}{+}\discretionary{}{\Wrappedafterbreak}{}\PY{n}{s12}\PY{o}{*}\PY{n}{s13}\PY{o}{+}\discretionary{}{\Wrappedafterbreak}{}\PY{n}{s13}\PY{o}{*}\PY{n}{s14}\PY{o}{+}\discretionary{}{\Wrappedafterbreak}{}\PY{n}{s14}\PY{o}{*}\PY{n}{s15}\PY{o}{+}\discretionary{}{\Wrappedafterbreak}{}\PY{n}{s15}\PY{o}{*}\PY{n}{s16}\PY{o}{+}\discretionary{}{\Wrappedafterbreak}{}\PY{n}{s16}\PY{o}{*}\PY{n}{s17}\PY{o}{+}\discretionary{}{\Wrappedafterbreak}{}\PY{n}{s17}\PY{o}{*}\PY{n}{s18}\PY{o}{+}\discretionary{}{\Wrappedafterbreak}{}\PY{n}{s18}\PY{o}{*}\PY{n}{s19}\PY{o}{+}\discretionary{}{\Wrappedafterbreak}{}\PY{n}{s19}\PY{o}{*}\PY{n}{s20}\PY{p}{)}
                    \PY{n}{M\PYZus{}20}\PY{o}{=}\PY{p}{(}\PY{n}{s1}\PY{o}{+}\PY{n}{s2}\PY{o}{+}\PY{n}{s3}\PY{o}{+}\PY{n}{s4}\PY{o}{+}\PY{n}{s5}\PY{o}{+}\PY{n}{s6}\PY{o}{+}\PY{n}{s7}\PY{o}{+}\PY{n}{s8}\PY{o}{+}\PY{n}{s9}\PY{o}{+}\PY{n}{s10}\PY{o}{+}\PY{n}{s11}\PY{o}{+}\PY{n}{s12}\PY{o}{+}\PY{n}{s13}\PY{o}{+}\PY{n}{s14}\PY{o}{+}\PY{n}{s15}\PY{o}{+}\PY{n}{s16}\PY{o}{+}\PY{n}{s17}\PY{o}{+}\PY{n}{s18}\PY{o}{+}\PY{n}{s19}\PY{o}{+}\PY{n}{s20}\PY{p}{)}

                    \PY{k}{if} \PY{p}{(}\PY{n}{E\PYZus{}20}\PY{p}{,} \PY{n}{M\PYZus{}20}\PY{p}{)} \PY{o+ow}{in} \PY{n}{distr}\PY{p}{:}
                     \PY{n}{distr}\PY{p}{[}\PY{p}{(}\PY{n}{E\PYZus{}20}\PY{p}{,} \PY{n}{M\PYZus{}20}\PY{p}{)}\PY{p}{]}\PY{o}{+}\PY{o}{=}\PY{l+m+mi}{1}
                    \PY{k}{else}\PY{p}{:}
                     \PY{n}{distr}\PY{o}{.}\PY{n}{update}\PY{p}{(}\PY{p}{\PYZob{}}\PY{p}{(}\PY{n}{E\PYZus{}20}\PY{p}{,} \PY{n}{M\PYZus{}20}\PY{p}{)} \PY{p}{:} \PY{l+m+mi}{1}\PY{p}{\PYZcb{}}\PY{p}{)}

\PY{n+nb}{print}\PY{p}{(}\PY{l+s+s2}{\PYZdq{}}\PY{l+s+s2}{SUM\PYZus{}E,M OMEGA(E,M) =}\PY{l+s+s2}{\PYZdq{}}\PY{p}{,} \PY{n+nb}{sum}\PY{p}{(}\PY{n}{distr}\PY{o}{.}\PY{n}{values}\PY{p}{(}\PY{p}{)}\PY{p}{)}\PY{p}{)}
\PY{k}{if} \PY{n+nb}{sum}\PY{p}{(}\PY{n}{distr}\PY{o}{.}\PY{n}{values}\PY{p}{(}\PY{p}{)}\PY{p}{)} \PY{o}{==} \PY{l+m+mi}{1048576}\PY{p}{:}
 \PY{n+nb}{print}\PY{p}{(}\PY{l+s+s2}{\PYZdq{}}\PY{l+s+s2}{\PYZhy{}\PYZgt{} check passed}\PY{l+s+s2}{\PYZdq{}}\PY{p}{)}
\end{Verbatim}
\end{tcolorbox}

    \begin{Verbatim}[commandchars=\\\{\}]
SUM\_E,M OMEGA(E,M) = 1048576
-> check passed
    \end{Verbatim}

    With \(\Omega(E,M)\) the exact values for \(N=20\) are, as in eq. (8)
and (9) of Homework 3,

\[
\langle A\rangle=\frac{\sum_{E,M}A\,\Omega(E,M)e^{-\beta E}}{\sum_{E,M}\Omega(E,M)e^{-\beta E}} \tag{12}
\]

for \(A=E, E^2, |M|, M^2\).

    \begin{tcolorbox}[breakable, size=fbox, boxrule=1pt, pad at break*=1mm,colback=cellbackground, colframe=cellborder]
\prompt{In}{incolor}{20}{\boxspacing}
\begin{Verbatim}[commandchars=\\\{\}]
\PY{k}{def} \PY{n+nf}{exact\PYZus{}enumeration}\PY{p}{(}\PY{n}{beta}\PY{p}{)}\PY{p}{:}
    \PY{n}{N} \PY{o}{=} \PY{l+m+mi}{20}
    \PY{n}{Z} \PY{o}{=} \PY{l+m+mi}{0}
    \PY{n}{E\PYZus{}1} \PY{o}{=} \PY{l+m+mi}{0}
    \PY{n}{E\PYZus{}2} \PY{o}{=} \PY{l+m+mi}{0}
    \PY{n}{M\PYZus{}abs} \PY{o}{=} \PY{l+m+mi}{0}
    \PY{n}{M\PYZus{}2} \PY{o}{=} \PY{l+m+mi}{0}

    \PY{k}{for} \PY{n}{E}\PY{p}{,} \PY{n}{M} \PY{o+ow}{in} \PY{n}{distr}\PY{p}{:}
        \PY{n}{w} \PY{o}{=} \PY{n}{distr}\PY{p}{[}\PY{n}{E}\PY{p}{,} \PY{n}{M}\PY{p}{]}\PY{o}{*}\PY{n}{exp}\PY{p}{(}\PY{o}{\PYZhy{}}\PY{n}{beta}\PY{o}{*}\PY{n}{E}\PY{p}{)}
        \PY{n}{Z} \PY{o}{+}\PY{o}{=} \PY{n}{w}
        \PY{n}{E\PYZus{}1} \PY{o}{+}\PY{o}{=} \PY{n}{E}\PY{o}{*}\PY{n}{w}
        \PY{n}{E\PYZus{}2} \PY{o}{+}\PY{o}{=} \PY{n}{E}\PY{o}{*}\PY{n}{E}\PY{o}{*}\PY{n}{w}
        \PY{n}{M\PYZus{}abs} \PY{o}{+}\PY{o}{=} \PY{n+nb}{abs}\PY{p}{(}\PY{n}{M}\PY{p}{)}\PY{o}{*}\PY{n}{w}
        \PY{n}{M\PYZus{}2} \PY{o}{+}\PY{o}{=} \PY{n}{M}\PY{o}{*}\PY{n}{M}\PY{o}{*}\PY{n}{w}

    \PY{n}{E\PYZus{}1}\PY{p}{,} \PY{n}{E\PYZus{}2}\PY{p}{,} \PY{n}{M\PYZus{}abs}\PY{p}{,} \PY{n}{M\PYZus{}2} \PY{o}{=} \PY{n}{E\PYZus{}1}\PY{o}{/}\PY{n}{Z}\PY{p}{,} \PY{n}{E\PYZus{}2}\PY{o}{/}\PY{n}{Z}\PY{p}{,} \PY{n}{M\PYZus{}abs}\PY{o}{/}\PY{n}{Z}\PY{p}{,} \PY{n}{M\PYZus{}2}\PY{o}{/}\PY{n}{Z}
    \PY{k}{return} \PY{n}{E\PYZus{}1}\PY{o}{/}\PY{n}{N}\PY{p}{,} \PY{n}{beta}\PY{o}{*}\PY{o}{*}\PY{l+m+mi}{2}\PY{o}{*}\PY{p}{(}\PY{n}{E\PYZus{}2} \PY{o}{\PYZhy{}} \PY{n}{E\PYZus{}1}\PY{o}{*}\PY{o}{*}\PY{l+m+mi}{2}\PY{p}{)}\PY{o}{/}\PY{n}{N}\PY{p}{,} \PY{n}{M\PYZus{}abs}\PY{o}{/}\PY{n}{N}\PY{p}{,} \PY{n}{beta}\PY{o}{*}\PY{n}{M\PYZus{}2}\PY{o}{/}\PY{n}{N}

\PY{n}{e\PYZus{}enum} \PY{o}{=} \PY{n}{np}\PY{o}{.}\PY{n}{array}\PY{p}{(}\PY{p}{[}\PY{n}{exact\PYZus{}enumeration}\PY{p}{(}\PY{n}{beta}\PY{p}{)}\PY{p}{[}\PY{l+m+mi}{0}\PY{p}{]} \PY{k}{for} \PY{n}{beta} \PY{o+ow}{in} \PY{n}{betas}\PY{p}{]}\PY{p}{)}
\PY{n}{c\PYZus{}enum} \PY{o}{=} \PY{n}{np}\PY{o}{.}\PY{n}{array}\PY{p}{(}\PY{p}{[}\PY{n}{exact\PYZus{}enumeration}\PY{p}{(}\PY{n}{beta}\PY{p}{)}\PY{p}{[}\PY{l+m+mi}{1}\PY{p}{]} \PY{k}{for} \PY{n}{beta} \PY{o+ow}{in} \PY{n}{betas}\PY{p}{]}\PY{p}{)}
\PY{n}{m\PYZus{}enum} \PY{o}{=} \PY{n}{np}\PY{o}{.}\PY{n}{array}\PY{p}{(}\PY{p}{[}\PY{n}{exact\PYZus{}enumeration}\PY{p}{(}\PY{n}{beta}\PY{p}{)}\PY{p}{[}\PY{l+m+mi}{2}\PY{p}{]} \PY{k}{for} \PY{n}{beta} \PY{o+ow}{in} \PY{n}{betas}\PY{p}{]}\PY{p}{)}
\PY{n}{chi\PYZus{}enum} \PY{o}{=} \PY{n}{np}\PY{o}{.}\PY{n}{array}\PY{p}{(}\PY{p}{[}\PY{n}{exact\PYZus{}enumeration}\PY{p}{(}\PY{n}{beta}\PY{p}{)}\PY{p}{[}\PY{l+m+mi}{3}\PY{p}{]} \PY{k}{for} \PY{n}{beta} \PY{o+ow}{in} \PY{n}{betas}\PY{p}{]}\PY{p}{)}
\end{Verbatim}
\end{tcolorbox}

    Comparison for \(N=20\). First we compare the numbers at a few values of
\(\beta\).

    \begin{tcolorbox}[breakable, size=fbox, boxrule=1pt, pad at break*=1mm,colback=cellbackground, colframe=cellborder]
\prompt{In}{incolor}{21}{\boxspacing}
\begin{Verbatim}[commandchars=\\\{\}]
\PY{n+nb}{print}\PY{p}{(}\PY{l+s+s2}{\PYZdq{}}\PY{l+s+s2}{COMPARISON FOR N=20 (Monte Carlo / exact enumeration / analytical):}\PY{l+s+s2}{\PYZdq{}}\PY{p}{)}
\PY{k}{for} \PY{n}{k} \PY{o+ow}{in} \PY{p}{[}\PY{l+m+mi}{5}\PY{p}{,} \PY{l+m+mi}{10}\PY{p}{,} \PY{l+m+mi}{15}\PY{p}{,} \PY{l+m+mi}{20}\PY{p}{]}\PY{p}{:}
    \PY{n}{beta} \PY{o}{=} \PY{n}{betas}\PY{p}{[}\PY{n}{k}\PY{p}{]}
    \PY{n+nb}{print}\PY{p}{(}\PY{l+s+sa}{f}\PY{l+s+s2}{\PYZdq{}}\PY{l+s+s2}{beta=}\PY{l+s+si}{\PYZob{}}\PY{n}{beta}\PY{l+s+si}{:}\PY{l+s+s2}{.1f}\PY{l+s+si}{\PYZcb{}}\PY{l+s+s2}{:}\PY{l+s+s2}{\PYZdq{}}\PY{p}{)}
    \PY{n+nb}{print}\PY{p}{(}\PY{l+s+sa}{f}\PY{l+s+s2}{\PYZdq{}}\PY{l+s+s2}{   e   = }\PY{l+s+si}{\PYZob{}}\PY{n}{e\PYZus{}MC}\PY{p}{[}\PY{l+m+mi}{20}\PY{p}{]}\PY{p}{[}\PY{n}{k}\PY{p}{]}\PY{l+s+si}{:}\PY{l+s+s2}{.4f}\PY{l+s+si}{\PYZcb{}}\PY{l+s+s2}{ / }\PY{l+s+si}{\PYZob{}}\PY{n}{e\PYZus{}enum}\PY{p}{[}\PY{n}{k}\PY{p}{]}\PY{l+s+si}{:}\PY{l+s+s2}{.4f}\PY{l+s+si}{\PYZcb{}}\PY{l+s+s2}{ / }\PY{l+s+si}{\PYZob{}}\PY{n}{e\PYZus{}analytic}\PY{p}{(}\PY{l+m+mi}{20}\PY{p}{,}\PY{+w}{ }\PY{n}{beta}\PY{p}{)}\PY{l+s+si}{:}\PY{l+s+s2}{.4f}\PY{l+s+si}{\PYZcb{}}\PY{l+s+s2}{\PYZdq{}}\PY{p}{)}
    \PY{n+nb}{print}\PY{p}{(}\PY{l+s+sa}{f}\PY{l+s+s2}{\PYZdq{}}\PY{l+s+s2}{   c   = }\PY{l+s+si}{\PYZob{}}\PY{n}{c\PYZus{}MC}\PY{p}{[}\PY{l+m+mi}{20}\PY{p}{]}\PY{p}{[}\PY{n}{k}\PY{p}{]}\PY{l+s+si}{:}\PY{l+s+s2}{.4f}\PY{l+s+si}{\PYZcb{}}\PY{l+s+s2}{ / }\PY{l+s+si}{\PYZob{}}\PY{n}{c\PYZus{}enum}\PY{p}{[}\PY{n}{k}\PY{p}{]}\PY{l+s+si}{:}\PY{l+s+s2}{.4f}\PY{l+s+si}{\PYZcb{}}\PY{l+s+s2}{ / }\PY{l+s+si}{\PYZob{}}\PY{n}{c\PYZus{}analytic}\PY{p}{(}\PY{l+m+mi}{20}\PY{p}{,}\PY{+w}{ }\PY{n}{beta}\PY{p}{)}\PY{l+s+si}{:}\PY{l+s+s2}{.4f}\PY{l+s+si}{\PYZcb{}}\PY{l+s+s2}{\PYZdq{}}\PY{p}{)}
    \PY{n+nb}{print}\PY{p}{(}\PY{l+s+sa}{f}\PY{l+s+s2}{\PYZdq{}}\PY{l+s+s2}{   m   = }\PY{l+s+si}{\PYZob{}}\PY{n}{m\PYZus{}MC}\PY{p}{[}\PY{l+m+mi}{20}\PY{p}{]}\PY{p}{[}\PY{n}{k}\PY{p}{]}\PY{l+s+si}{:}\PY{l+s+s2}{.4f}\PY{l+s+si}{\PYZcb{}}\PY{l+s+s2}{ / }\PY{l+s+si}{\PYZob{}}\PY{n}{m\PYZus{}enum}\PY{p}{[}\PY{n}{k}\PY{p}{]}\PY{l+s+si}{:}\PY{l+s+s2}{.4f}\PY{l+s+si}{\PYZcb{}}\PY{l+s+s2}{\PYZdq{}}\PY{p}{)}
    \PY{n+nb}{print}\PY{p}{(}\PY{l+s+sa}{f}\PY{l+s+s2}{\PYZdq{}}\PY{l+s+s2}{   chi = }\PY{l+s+si}{\PYZob{}}\PY{n}{chi\PYZus{}MC}\PY{p}{[}\PY{l+m+mi}{20}\PY{p}{]}\PY{p}{[}\PY{n}{k}\PY{p}{]}\PY{l+s+si}{:}\PY{l+s+s2}{.4f}\PY{l+s+si}{\PYZcb{}}\PY{l+s+s2}{ / }\PY{l+s+si}{\PYZob{}}\PY{n}{chi\PYZus{}enum}\PY{p}{[}\PY{n}{k}\PY{p}{]}\PY{l+s+si}{:}\PY{l+s+s2}{.4f}\PY{l+s+si}{\PYZcb{}}\PY{l+s+s2}{\PYZdq{}}\PY{p}{)}
\end{Verbatim}
\end{tcolorbox}

    \begin{Verbatim}[commandchars=\\\{\}]
COMPARISON FOR N=20 (Monte Carlo / exact enumeration / analytical):
beta=0.5:
   e   = -0.4392 / -0.4390 / -0.4390
   c   = 0.1872 / 0.1868 / 0.1868
   m   = 0.2878 / 0.2884
   chi = 1.2747 / 1.2793
beta=1.0:
   e   = -0.7224 / -0.7235 / -0.7235
   c   = 0.3974 / 0.3990 / 0.3990
   m   = 0.4574 / 0.4601
   chi = 5.9914 / 6.0549
beta=1.5:
   e   = -0.8594 / -0.8599 / -0.8599
   c   = 0.3858 / 0.3863 / 0.3863
   m   = 0.6726 / 0.6745
   chi = 17.0112 / 17.0937
beta=2.0:
   e   = -0.9149 / -0.9158 / -0.9158
   c   = 0.2748 / 0.2685 / 0.2685
   m   = 0.8412 / 0.8459
   chi = 31.5732 / 31.8077
    \end{Verbatim}

    Then we plot the Monte Carlo results (points) together with the exact
results (lines) for all \(\beta\).

    \begin{tcolorbox}[breakable, size=fbox, boxrule=1pt, pad at break*=1mm,colback=cellbackground, colframe=cellborder]
\prompt{In}{incolor}{22}{\boxspacing}
\begin{Verbatim}[commandchars=\\\{\}]
\PY{n}{beta\PYZus{}fine} \PY{o}{=} \PY{n}{np}\PY{o}{.}\PY{n}{linspace}\PY{p}{(}\PY{l+m+mi}{0}\PY{p}{,} \PY{l+m+mi}{2}\PY{p}{,} \PY{l+m+mi}{200}\PY{p}{)}

\PY{n}{plt}\PY{o}{.}\PY{n}{plot}\PY{p}{(}\PY{n}{betas}\PY{p}{,} \PY{n}{e\PYZus{}MC}\PY{p}{[}\PY{l+m+mi}{20}\PY{p}{]}\PY{p}{,} \PY{l+s+s2}{\PYZdq{}}\PY{l+s+s2}{o}\PY{l+s+s2}{\PYZdq{}}\PY{p}{,} \PY{n}{label}\PY{o}{=}\PY{l+s+s2}{\PYZdq{}}\PY{l+s+s2}{e, Monte Carlo}\PY{l+s+s2}{\PYZdq{}}\PY{p}{)}
\PY{n}{plt}\PY{o}{.}\PY{n}{plot}\PY{p}{(}\PY{n}{beta\PYZus{}fine}\PY{p}{,} \PY{n}{e\PYZus{}analytic}\PY{p}{(}\PY{l+m+mi}{20}\PY{p}{,} \PY{n}{beta\PYZus{}fine}\PY{p}{)}\PY{p}{,} \PY{n}{label}\PY{o}{=}\PY{l+s+s2}{\PYZdq{}}\PY{l+s+s2}{e, analytical}\PY{l+s+s2}{\PYZdq{}}\PY{p}{)}
\PY{n}{plt}\PY{o}{.}\PY{n}{plot}\PY{p}{(}\PY{n}{betas}\PY{p}{,} \PY{n}{c\PYZus{}MC}\PY{p}{[}\PY{l+m+mi}{20}\PY{p}{]}\PY{p}{,} \PY{l+s+s2}{\PYZdq{}}\PY{l+s+s2}{o}\PY{l+s+s2}{\PYZdq{}}\PY{p}{,} \PY{n}{label}\PY{o}{=}\PY{l+s+s2}{\PYZdq{}}\PY{l+s+s2}{c, Monte Carlo}\PY{l+s+s2}{\PYZdq{}}\PY{p}{)}
\PY{n}{plt}\PY{o}{.}\PY{n}{plot}\PY{p}{(}\PY{n}{beta\PYZus{}fine}\PY{p}{,} \PY{n}{c\PYZus{}analytic}\PY{p}{(}\PY{l+m+mi}{20}\PY{p}{,} \PY{n}{beta\PYZus{}fine}\PY{p}{)}\PY{p}{,} \PY{n}{label}\PY{o}{=}\PY{l+s+s2}{\PYZdq{}}\PY{l+s+s2}{c, analytical}\PY{l+s+s2}{\PYZdq{}}\PY{p}{)}
\PY{n}{plt}\PY{o}{.}\PY{n}{xlabel}\PY{p}{(}\PY{l+s+sa}{r}\PY{l+s+s2}{\PYZdq{}}\PY{l+s+s2}{\PYZdl{}}\PY{l+s+s2}{\PYZbs{}}\PY{l+s+s2}{beta\PYZdl{}}\PY{l+s+s2}{\PYZdq{}}\PY{p}{)}
\PY{n}{plt}\PY{o}{.}\PY{n}{legend}\PY{p}{(}\PY{p}{)}
\PY{n}{plt}\PY{o}{.}\PY{n}{show}\PY{p}{(}\PY{p}{)}
\end{Verbatim}
\end{tcolorbox}

    \begin{center}
    \adjustimage{max size={0.9\linewidth}{0.9\paperheight}}{CompSim_HW4_files/CompSim_HW4_36_0.png}
    \end{center}
    { \hspace*{\fill} \\}
    
    Figure 5. Mean energy \(e\) and specific heat \(c\) per spin as a
function of \(\beta\) for \(N=20\), Monte Carlo and analytical result of
Problem 1

    \begin{tcolorbox}[breakable, size=fbox, boxrule=1pt, pad at break*=1mm,colback=cellbackground, colframe=cellborder]
\prompt{In}{incolor}{23}{\boxspacing}
\begin{Verbatim}[commandchars=\\\{\}]
\PY{n}{plt}\PY{o}{.}\PY{n}{plot}\PY{p}{(}\PY{n}{betas}\PY{p}{,} \PY{n}{m\PYZus{}MC}\PY{p}{[}\PY{l+m+mi}{20}\PY{p}{]}\PY{p}{,} \PY{l+s+s2}{\PYZdq{}}\PY{l+s+s2}{o}\PY{l+s+s2}{\PYZdq{}}\PY{p}{,} \PY{n}{label}\PY{o}{=}\PY{l+s+s2}{\PYZdq{}}\PY{l+s+s2}{m, Monte Carlo}\PY{l+s+s2}{\PYZdq{}}\PY{p}{)}
\PY{n}{plt}\PY{o}{.}\PY{n}{plot}\PY{p}{(}\PY{n}{betas}\PY{p}{,} \PY{n}{m\PYZus{}enum}\PY{p}{,} \PY{n}{label}\PY{o}{=}\PY{l+s+s2}{\PYZdq{}}\PY{l+s+s2}{m, exact enumeration}\PY{l+s+s2}{\PYZdq{}}\PY{p}{)}
\PY{n}{plt}\PY{o}{.}\PY{n}{plot}\PY{p}{(}\PY{n}{betas}\PY{p}{,} \PY{n}{chi\PYZus{}MC}\PY{p}{[}\PY{l+m+mi}{20}\PY{p}{]}\PY{p}{,} \PY{l+s+s2}{\PYZdq{}}\PY{l+s+s2}{o}\PY{l+s+s2}{\PYZdq{}}\PY{p}{,} \PY{n}{label}\PY{o}{=}\PY{l+s+sa}{r}\PY{l+s+s2}{\PYZdq{}}\PY{l+s+s2}{\PYZdl{}}\PY{l+s+s2}{\PYZbs{}}\PY{l+s+s2}{chi\PYZdl{}, Monte Carlo}\PY{l+s+s2}{\PYZdq{}}\PY{p}{)}
\PY{n}{plt}\PY{o}{.}\PY{n}{plot}\PY{p}{(}\PY{n}{betas}\PY{p}{,} \PY{n}{chi\PYZus{}enum}\PY{p}{,} \PY{n}{label}\PY{o}{=}\PY{l+s+sa}{r}\PY{l+s+s2}{\PYZdq{}}\PY{l+s+s2}{\PYZdl{}}\PY{l+s+s2}{\PYZbs{}}\PY{l+s+s2}{chi\PYZdl{}, exact enumeration}\PY{l+s+s2}{\PYZdq{}}\PY{p}{)}
\PY{n}{plt}\PY{o}{.}\PY{n}{yscale}\PY{p}{(}\PY{l+s+s2}{\PYZdq{}}\PY{l+s+s2}{log}\PY{l+s+s2}{\PYZdq{}}\PY{p}{)}
\PY{n}{plt}\PY{o}{.}\PY{n}{xlabel}\PY{p}{(}\PY{l+s+sa}{r}\PY{l+s+s2}{\PYZdq{}}\PY{l+s+s2}{\PYZdl{}}\PY{l+s+s2}{\PYZbs{}}\PY{l+s+s2}{beta\PYZdl{}}\PY{l+s+s2}{\PYZdq{}}\PY{p}{)}
\PY{n}{plt}\PY{o}{.}\PY{n}{legend}\PY{p}{(}\PY{p}{)}
\PY{n}{plt}\PY{o}{.}\PY{n}{show}\PY{p}{(}\PY{p}{)}
\end{Verbatim}
\end{tcolorbox}

    \begin{center}
    \adjustimage{max size={0.9\linewidth}{0.9\paperheight}}{CompSim_HW4_files/CompSim_HW4_38_0.png}
    \end{center}
    { \hspace*{\fill} \\}
    
    Figure 6. Magnetisation \(m\) and susceptibility \(\chi\) per spin as a
function of \(\beta\) for \(N=20\), Monte Carlo and exact enumeration
(logarithmic y-scale)

    As we can see on Fig. 5 and Fig. 6, the Monte Carlo results for \(N=20\)
agree with the exact enumeration and the analytical results within about
\(1\%\).

The only exception is the magnetisation at \(\beta=0\). There every flip
is accepted, so one sweep is exactly \(N=20\) flips and \(M\) changes
only by multiples of \(4\) between measurements. The simulation
therefore visits only half of the configurations.

    Now we plot the results for all \(N\). For \(e\) and \(c\) we also plot
the analytical result eq. (11) as lines.

    \begin{tcolorbox}[breakable, size=fbox, boxrule=1pt, pad at break*=1mm,colback=cellbackground, colframe=cellborder]
\prompt{In}{incolor}{24}{\boxspacing}
\begin{Verbatim}[commandchars=\\\{\}]
\PY{n}{colors} \PY{o}{=} \PY{p}{[}\PY{l+s+s2}{\PYZdq{}}\PY{l+s+s2}{C0}\PY{l+s+s2}{\PYZdq{}}\PY{p}{,} \PY{l+s+s2}{\PYZdq{}}\PY{l+s+s2}{C1}\PY{l+s+s2}{\PYZdq{}}\PY{p}{,} \PY{l+s+s2}{\PYZdq{}}\PY{l+s+s2}{C2}\PY{l+s+s2}{\PYZdq{}}\PY{p}{,} \PY{l+s+s2}{\PYZdq{}}\PY{l+s+s2}{C3}\PY{l+s+s2}{\PYZdq{}}\PY{p}{]}

\PY{k}{for} \PY{n}{N}\PY{p}{,} \PY{n}{col} \PY{o+ow}{in} \PY{n+nb}{zip}\PY{p}{(}\PY{n}{N\PYZus{}values}\PY{p}{,} \PY{n}{colors}\PY{p}{)}\PY{p}{:}
    \PY{n}{plt}\PY{o}{.}\PY{n}{plot}\PY{p}{(}\PY{n}{betas}\PY{p}{,} \PY{n}{e\PYZus{}MC}\PY{p}{[}\PY{n}{N}\PY{p}{]}\PY{p}{,} \PY{l+s+s2}{\PYZdq{}}\PY{l+s+s2}{o}\PY{l+s+s2}{\PYZdq{}}\PY{p}{,} \PY{n}{color}\PY{o}{=}\PY{n}{col}\PY{p}{,} \PY{n}{label}\PY{o}{=}\PY{l+s+sa}{f}\PY{l+s+s2}{\PYZdq{}}\PY{l+s+s2}{N=}\PY{l+s+si}{\PYZob{}}\PY{n}{N}\PY{l+s+si}{\PYZcb{}}\PY{l+s+s2}{\PYZdq{}}\PY{p}{)}
    \PY{n}{plt}\PY{o}{.}\PY{n}{plot}\PY{p}{(}\PY{n}{beta\PYZus{}fine}\PY{p}{,} \PY{n}{e\PYZus{}analytic}\PY{p}{(}\PY{n}{N}\PY{p}{,} \PY{n}{beta\PYZus{}fine}\PY{p}{)}\PY{p}{,} \PY{n}{color}\PY{o}{=}\PY{n}{col}\PY{p}{)}
\PY{n}{plt}\PY{o}{.}\PY{n}{xlabel}\PY{p}{(}\PY{l+s+sa}{r}\PY{l+s+s2}{\PYZdq{}}\PY{l+s+s2}{\PYZdl{}}\PY{l+s+s2}{\PYZbs{}}\PY{l+s+s2}{beta\PYZdl{}}\PY{l+s+s2}{\PYZdq{}}\PY{p}{)}
\PY{n}{plt}\PY{o}{.}\PY{n}{ylabel}\PY{p}{(}\PY{l+s+s2}{\PYZdq{}}\PY{l+s+s2}{e}\PY{l+s+s2}{\PYZdq{}}\PY{p}{)}
\PY{n}{plt}\PY{o}{.}\PY{n}{legend}\PY{p}{(}\PY{p}{)}
\PY{n}{plt}\PY{o}{.}\PY{n}{show}\PY{p}{(}\PY{p}{)}
\end{Verbatim}
\end{tcolorbox}

    \begin{center}
    \adjustimage{max size={0.9\linewidth}{0.9\paperheight}}{CompSim_HW4_files/CompSim_HW4_42_0.png}
    \end{center}
    { \hspace*{\fill} \\}
    
    Figure 7. Mean energy per spin \(e\) as a function of \(\beta\) for
\(N=20,100,1000,10000\). Points: Monte Carlo, lines: analytical result

    \begin{tcolorbox}[breakable, size=fbox, boxrule=1pt, pad at break*=1mm,colback=cellbackground, colframe=cellborder]
\prompt{In}{incolor}{25}{\boxspacing}
\begin{Verbatim}[commandchars=\\\{\}]
\PY{k}{for} \PY{n}{N}\PY{p}{,} \PY{n}{col} \PY{o+ow}{in} \PY{n+nb}{zip}\PY{p}{(}\PY{n}{N\PYZus{}values}\PY{p}{,} \PY{n}{colors}\PY{p}{)}\PY{p}{:}
    \PY{n}{plt}\PY{o}{.}\PY{n}{plot}\PY{p}{(}\PY{n}{betas}\PY{p}{,} \PY{n}{c\PYZus{}MC}\PY{p}{[}\PY{n}{N}\PY{p}{]}\PY{p}{,} \PY{l+s+s2}{\PYZdq{}}\PY{l+s+s2}{o}\PY{l+s+s2}{\PYZdq{}}\PY{p}{,} \PY{n}{color}\PY{o}{=}\PY{n}{col}\PY{p}{,} \PY{n}{label}\PY{o}{=}\PY{l+s+sa}{f}\PY{l+s+s2}{\PYZdq{}}\PY{l+s+s2}{N=}\PY{l+s+si}{\PYZob{}}\PY{n}{N}\PY{l+s+si}{\PYZcb{}}\PY{l+s+s2}{\PYZdq{}}\PY{p}{)}
    \PY{n}{plt}\PY{o}{.}\PY{n}{plot}\PY{p}{(}\PY{n}{beta\PYZus{}fine}\PY{p}{,} \PY{n}{c\PYZus{}analytic}\PY{p}{(}\PY{n}{N}\PY{p}{,} \PY{n}{beta\PYZus{}fine}\PY{p}{)}\PY{p}{,} \PY{n}{color}\PY{o}{=}\PY{n}{col}\PY{p}{)}
\PY{n}{plt}\PY{o}{.}\PY{n}{xlabel}\PY{p}{(}\PY{l+s+sa}{r}\PY{l+s+s2}{\PYZdq{}}\PY{l+s+s2}{\PYZdl{}}\PY{l+s+s2}{\PYZbs{}}\PY{l+s+s2}{beta\PYZdl{}}\PY{l+s+s2}{\PYZdq{}}\PY{p}{)}
\PY{n}{plt}\PY{o}{.}\PY{n}{ylabel}\PY{p}{(}\PY{l+s+s2}{\PYZdq{}}\PY{l+s+s2}{c}\PY{l+s+s2}{\PYZdq{}}\PY{p}{)}
\PY{n}{plt}\PY{o}{.}\PY{n}{legend}\PY{p}{(}\PY{p}{)}
\PY{n}{plt}\PY{o}{.}\PY{n}{show}\PY{p}{(}\PY{p}{)}
\end{Verbatim}
\end{tcolorbox}

    \begin{center}
    \adjustimage{max size={0.9\linewidth}{0.9\paperheight}}{CompSim_HW4_files/CompSim_HW4_44_0.png}
    \end{center}
    { \hspace*{\fill} \\}
    
    Figure 8. Specific heat per spin \(c\) as a function of \(\beta\) for
\(N=20,100,1000,10000\). Points: Monte Carlo, lines: analytical result

    \begin{tcolorbox}[breakable, size=fbox, boxrule=1pt, pad at break*=1mm,colback=cellbackground, colframe=cellborder]
\prompt{In}{incolor}{26}{\boxspacing}
\begin{Verbatim}[commandchars=\\\{\}]
\PY{k}{for} \PY{n}{N} \PY{o+ow}{in} \PY{n}{N\PYZus{}values}\PY{p}{:}
    \PY{n}{plt}\PY{o}{.}\PY{n}{plot}\PY{p}{(}\PY{n}{betas}\PY{p}{,} \PY{n}{m\PYZus{}MC}\PY{p}{[}\PY{n}{N}\PY{p}{]}\PY{p}{,} \PY{l+s+s2}{\PYZdq{}}\PY{l+s+s2}{o\PYZhy{}}\PY{l+s+s2}{\PYZdq{}}\PY{p}{,} \PY{n}{label}\PY{o}{=}\PY{l+s+sa}{f}\PY{l+s+s2}{\PYZdq{}}\PY{l+s+s2}{N=}\PY{l+s+si}{\PYZob{}}\PY{n}{N}\PY{l+s+si}{\PYZcb{}}\PY{l+s+s2}{\PYZdq{}}\PY{p}{)}
\PY{n}{plt}\PY{o}{.}\PY{n}{xlabel}\PY{p}{(}\PY{l+s+sa}{r}\PY{l+s+s2}{\PYZdq{}}\PY{l+s+s2}{\PYZdl{}}\PY{l+s+s2}{\PYZbs{}}\PY{l+s+s2}{beta\PYZdl{}}\PY{l+s+s2}{\PYZdq{}}\PY{p}{)}
\PY{n}{plt}\PY{o}{.}\PY{n}{ylabel}\PY{p}{(}\PY{l+s+s2}{\PYZdq{}}\PY{l+s+s2}{m}\PY{l+s+s2}{\PYZdq{}}\PY{p}{)}
\PY{n}{plt}\PY{o}{.}\PY{n}{legend}\PY{p}{(}\PY{p}{)}
\PY{n}{plt}\PY{o}{.}\PY{n}{show}\PY{p}{(}\PY{p}{)}
\end{Verbatim}
\end{tcolorbox}

    \begin{center}
    \adjustimage{max size={0.9\linewidth}{0.9\paperheight}}{CompSim_HW4_files/CompSim_HW4_46_0.png}
    \end{center}
    { \hspace*{\fill} \\}
    
    Figure 9. Magnetisation per spin \(m=\langle|M|\rangle/N\) as a function
of \(\beta\) for \(N=20,100,1000,10000\)

    \begin{tcolorbox}[breakable, size=fbox, boxrule=1pt, pad at break*=1mm,colback=cellbackground, colframe=cellborder]
\prompt{In}{incolor}{27}{\boxspacing}
\begin{Verbatim}[commandchars=\\\{\}]
\PY{k}{for} \PY{n}{N} \PY{o+ow}{in} \PY{n}{N\PYZus{}values}\PY{p}{:}
    \PY{n}{plt}\PY{o}{.}\PY{n}{plot}\PY{p}{(}\PY{n}{betas}\PY{p}{,} \PY{n}{chi\PYZus{}MC}\PY{p}{[}\PY{n}{N}\PY{p}{]}\PY{p}{,} \PY{l+s+s2}{\PYZdq{}}\PY{l+s+s2}{o\PYZhy{}}\PY{l+s+s2}{\PYZdq{}}\PY{p}{,} \PY{n}{label}\PY{o}{=}\PY{l+s+sa}{f}\PY{l+s+s2}{\PYZdq{}}\PY{l+s+s2}{N=}\PY{l+s+si}{\PYZob{}}\PY{n}{N}\PY{l+s+si}{\PYZcb{}}\PY{l+s+s2}{\PYZdq{}}\PY{p}{)}
\PY{n}{plt}\PY{o}{.}\PY{n}{yscale}\PY{p}{(}\PY{l+s+s2}{\PYZdq{}}\PY{l+s+s2}{log}\PY{l+s+s2}{\PYZdq{}}\PY{p}{)}
\PY{n}{plt}\PY{o}{.}\PY{n}{xlabel}\PY{p}{(}\PY{l+s+sa}{r}\PY{l+s+s2}{\PYZdq{}}\PY{l+s+s2}{\PYZdl{}}\PY{l+s+s2}{\PYZbs{}}\PY{l+s+s2}{beta\PYZdl{}}\PY{l+s+s2}{\PYZdq{}}\PY{p}{)}
\PY{n}{plt}\PY{o}{.}\PY{n}{ylabel}\PY{p}{(}\PY{l+s+sa}{r}\PY{l+s+s2}{\PYZdq{}}\PY{l+s+s2}{\PYZdl{}}\PY{l+s+s2}{\PYZbs{}}\PY{l+s+s2}{chi\PYZdl{}}\PY{l+s+s2}{\PYZdq{}}\PY{p}{)}
\PY{n}{plt}\PY{o}{.}\PY{n}{legend}\PY{p}{(}\PY{p}{)}
\PY{n}{plt}\PY{o}{.}\PY{n}{show}\PY{p}{(}\PY{p}{)}
\end{Verbatim}
\end{tcolorbox}

    \begin{center}
    \adjustimage{max size={0.9\linewidth}{0.9\paperheight}}{CompSim_HW4_files/CompSim_HW4_48_0.png}
    \end{center}
    { \hspace*{\fill} \\}
    
    Figure 10. Susceptibility per spin \(\chi\) as a function of \(\beta\)
for \(N=20,100,1000,10000\) (logarithmic y-scale)

    \begin{tcolorbox}[breakable, size=fbox, boxrule=1pt, pad at break*=1mm,colback=cellbackground, colframe=cellborder]
\prompt{In}{incolor}{28}{\boxspacing}
\begin{Verbatim}[commandchars=\\\{\}]
\PY{k}{for} \PY{n}{N} \PY{o+ow}{in} \PY{n}{N\PYZus{}values}\PY{p}{:}
    \PY{n}{plt}\PY{o}{.}\PY{n}{plot}\PY{p}{(}\PY{n}{betas}\PY{p}{,} \PY{n}{acc\PYZus{}MC}\PY{p}{[}\PY{n}{N}\PY{p}{]}\PY{p}{,} \PY{l+s+s2}{\PYZdq{}}\PY{l+s+s2}{o\PYZhy{}}\PY{l+s+s2}{\PYZdq{}}\PY{p}{,} \PY{n}{label}\PY{o}{=}\PY{l+s+sa}{f}\PY{l+s+s2}{\PYZdq{}}\PY{l+s+s2}{N=}\PY{l+s+si}{\PYZob{}}\PY{n}{N}\PY{l+s+si}{\PYZcb{}}\PY{l+s+s2}{\PYZdq{}}\PY{p}{)}
\PY{n}{plt}\PY{o}{.}\PY{n}{xlabel}\PY{p}{(}\PY{l+s+sa}{r}\PY{l+s+s2}{\PYZdq{}}\PY{l+s+s2}{\PYZdl{}}\PY{l+s+s2}{\PYZbs{}}\PY{l+s+s2}{beta\PYZdl{}}\PY{l+s+s2}{\PYZdq{}}\PY{p}{)}
\PY{n}{plt}\PY{o}{.}\PY{n}{ylabel}\PY{p}{(}\PY{l+s+s2}{\PYZdq{}}\PY{l+s+s2}{acceptance rate}\PY{l+s+s2}{\PYZdq{}}\PY{p}{)}
\PY{n}{plt}\PY{o}{.}\PY{n}{legend}\PY{p}{(}\PY{p}{)}
\PY{n}{plt}\PY{o}{.}\PY{n}{show}\PY{p}{(}\PY{p}{)}
\end{Verbatim}
\end{tcolorbox}

    \begin{center}
    \adjustimage{max size={0.9\linewidth}{0.9\paperheight}}{CompSim_HW4_files/CompSim_HW4_50_0.png}
    \end{center}
    { \hspace*{\fill} \\}
    
    Figure 11. Acceptance rate of the spin flips as a function of \(\beta\)
for \(N=20,100,1000,10000\)


    % Add a bibliography block to the postdoc
    
    
    
\end{document}
